\chapter{Analytic Control and Noncollapsing}

This chapter establishes the analytic toolkit for controlled Ricci flows with
surgery.  Reduced geometry, local curvature estimates, compactness, and
noncollapsing provide the inputs used to analyze singularity models and to
continue the flow in the next chapter.

\section{Surgery Flows and Fixed Controls}

\begin{definition}[Ricci flow with surgery]
\label{def:ricci-flow-with-surgery}
\dcref{MT:ch14:14.8}
\source{morgan-tian}
\notready
A \emph{Ricci flow with surgery} is a pair $(\mathcal M,G)$ with the
following structure.  The surgery spacetime $\mathcal M$ has a continuous
time function $\mathbf t\colon\mathcal M\to I$ which is $C^\infty$ with
nonvanishing differential on every ordinary or one-sided smooth chart, a
time vector field $\chi$ which is smooth on those charts and satisfies
$\chi(\mathbf t)=1$, and the ordinary, one-sided, and surgery charts
which distinguish smooth, exposed, and singular points.  Every regular slice
$M_t=\mathbf t^{-1}(t)$ is a smooth manifold, and the horizontal bundle is
$\mathcal H=\ker(d\mathbf t)$.  The tensor $G$ is a smooth metric on
$\mathcal H$ and satisfies
\[
 \mathcal L_\chi G=-2\operatorname{Ric}_G.
\]
At an exposed point or an endpoint of $I$, the Lie derivative is one-sided.
At each transition time $t=t_k^+=t_{k+1}^-$ there are a backward slice
$\mathcal M_t^- = \Omega_k\times\{t\}$, where $\Omega_k$ is the open set of
points with bounded curvature and smooth terminal limit on the preceding
piece, and a forward slice $\mathcal M_t^+=M_{k+1}\times\{t\}$, the initial
slice of the following piece; at regular times
$\mathcal M_t^- = \mathcal M_t^+=M_t$, while at a terminal singular time only
$\mathcal M_t^-$ is defined.  The path metrics on these slices may be
infinite between distinct components.  The flow is three-dimensional when
its slices are three-manifolds.

For $x$ in a slice side and $s>0$, let
$B^\pm(x,t,s)$ denote the open radius-$s$ ball in the path metric of the
connected component of $\mathcal M_t^\pm$ containing $x$; at a regular time the
two signs agree with the ordinary metric ball.  The metric on the backward
side is the terminal presurgery metric and the metric on the forward side is
the postsurgery initial metric.  Thus a ball or parabolic region with a minus
(respectively plus) sign is formed only from points on the corresponding side
of a transition.

An admissible curve is a continuous curve whose time-$u$ point lies in the
corresponding slice and whose restriction to each constituent flow is smooth.
It is \emph{static} when its spatial component is constant in every chart.
For $(x,t)\in\mathcal M_t^+$, with $t+\Delta t$ in the time domain and
$\Delta t>0$, let
$P(x,t,s,\Delta t)$ be the union of maximal static curves starting in
$B^+(x,t,s)$ and defined through $t+\Delta t$.  For
$(x,t)\in\mathcal M_t^-$ and $\Delta t<0$, define
$P(x,t,s,\Delta t)$ by maximal static curves ending in $B^-(x,t,s)$ and
defined back to $t+\Delta t$.  No forward development is formed from a
terminal singular slice, and no backward development is formed outside the
time domain.  A set $Y\subset\mathcal M_t$ is
\emph{unscathed on $[c,d]$} if every point of $Y$ lies on a static curve
throughout that interval.  A parabolic region is unscathed when its ball is
unscathed and the region is its full static development.
\end{definition}

\begin{definition}[Noncollapsing for a surgery flow]
\label{def:surgery-flow-noncollapsing}
\dcref{MT:ch14:14.14,MT:ch14:14.15}
\source{morgan-tian}
\uses{def:ricci-flow-with-surgery}
\notready
Let $(\mathcal M,G)$ be a three-dimensional Ricci flow with surgery.  Let
$\kappa:[0,\infty)\to(0,\infty)$ be a positive function and let
$\epsilon>0$.  The flow is \emph{$\kappa$-noncollapsed below $\epsilon$} if,
for every $(x,t)\in\mathcal M_t^-$ and
$0<\rho<\epsilon$ with $t\ge\rho^2$,
\[
 P(x,t,\rho,-\rho^2)\text{ is unscathed},\qquad
 |\operatorname{Rm}|\le \rho^{-2}\text{ on }P(x,t,\rho,-\rho^2)
\quad\Longrightarrow\quad
 \left[\operatorname{Vol}_{g(t)}B^-(x,t,\rho)\ge\kappa(t)\rho^3
 \ \text{or}\ x\text{ lies in a component of }M_t^-
 \text{ with positive sectional curvature}\right].
\]
This is the source convention: a compact positive-sectional-curvature
component is an allowed exceptional branch.  At a regular time this is the
same predicate with $B^-=B$; at a transition it is evaluated on the backward
slice.  Whenever a later argument needs an actual volume lower bound, it must
assume the nonpositive branch (or invoke a separate positive-component
estimate).
\end{definition}

\begin{definition}[Normalized compact Riemannian three-manifold]
\label{def:normalized-riemannian-three-manifold}
\dcref{KL:II.5:Def1}
\source{kleiner-lott}
\notready
A compact Riemannian three-manifold $(M,g)$ is normalized when
$|\operatorname{Rm}_g|\le1$ and
\[
 \operatorname{Vol}_g B_g(x,1)\ge
 \tfrac12\operatorname{Vol}_{\mathbb R^3}B(0,1)
 \qquad(x\in M).
\]
\end{definition}

\begin{lemma}[Imported interface: fixed standard surgery-cap model]
\label{lem:standard-surgery-cap-model}
\dcref{KL:II.2:before-60.2,KL:II.4:72.1,KL:II.4:72.20}
\source{kleiner-lott}
\notready
Fix a sufficiently small global $\epsilon>0$ and put $A=10$.  There are
$B>A$, $k>0$, a smooth $F:(-B,\infty)\to\mathbb R$, and a complete
$O(3)$-invariant metric $g_0$ on
$\mathbb R^3=D^3\cup((-B,\infty)\times S^2)$ such that $F=0$ on
$[0,\infty)$, $g_0=e^{2F}g_{\mathrm{cyl}}$ on the cylindrical end, and
$g_0|_{D^3}$ has constant sectional curvature $k^2$.  The restriction
$f=F|_{(-A,0]}$ satisfies
\[
 f^{(m)}(0)=0\ (m\ge0),\quad f<0,\quad f'>0,\quad f''<0,\quad
 \|f\|_{C^2}<\epsilon,\quad
 \max\{|f|,|f'|\}\le\epsilon|f''|,
\]
and $g_0$ has nonnegative curvature operator and scalar curvature at least
one.  Fix this tuple $(B,F,k,g_0)$ henceforth.
\end{lemma}

\begin{proof}
Choose a smooth concave profile on $(-A,0]$ with the prescribed flat jet at
$0$, negative value, and positive derivative, and scale it so that its $C^2$
norm is below $\epsilon$.  The warped-product curvature formulas for
$e^{2F}(ds^2+g_{S^2})$ show that the profile has nonnegative curvature
operator on the neck.  Extend it by the constant cylindrical profile for
$s\ge0$ and solve the radial cap ODE on $(-B,-A)$, choosing $B>A$ so that the
warping function closes smoothly.  The matching conditions give a round ball
of sectional curvature $k^2$ and preserve nonnegative curvature; increasing
$B$ if necessary gives scalar curvature at least one.  The resulting radial
metric is complete and $O(3)$-invariant.  This is the fixed cap construction
in Kleiner--Lott II.2 and II.4, with all constants chosen before any flow or
surgery time is introduced.
\end{proof}

\begin{definition}[Standard solution]
\label{def:standard-solution}
\dcref{KL:II.2:60.2}
\source{kleiner-lott}
\uses{lem:standard-surgery-cap-model}
\notready
The standard solution is the maximal Ricci flow $(\mathbb R^3,g(t))$ with
$g(0)=g_0$ and curvature bounded on every compact subinterval of its domain.
\end{definition}

\begin{definition}[Hamilton--Ivey pinching condition]
\label{def:hamilton-ivey-pinching-condition}
\dcref{KL:appB:B.5}
\source{kleiner-lott}
\notready
For $t>0$, let $\lambda_1\le\lambda_2\le\lambda_3$ be the curvature-operator
eigenvalues.  The time-$t$ Hamilton--Ivey condition is automatic when
$\lambda_1\ge0$ and otherwise, with $X=-\lambda_1$,
\[
 tR\ge tX\left(\log(tX)+\log\frac{1+t}{t}-3\right).
\]
\end{definition}

\begin{lemma}[Imported interface: Hamilton--Ivey envelope]
\label{lem:hamilton-ivey-phi-envelope}
\dcref{KL:appB:B.5,KL:appB:after-B.6,KL:appB:phi}
\uses{def:hamilton-ivey-pinching-condition}
\source{kleiner-lott}
\notready
For every finite horizon $\widehat T>0$, every complete Ricci flow
$(M,g(t))_{0\le t\le\widehat T}$ with curvature bounded on compact time
intervals, and every initial lower bound
$\lambda_1(\cdot,0)\ge-1$, there is a positive smooth function
$\Phi_{\mathrm{HI}}^{\widehat T,g(0)}$ on $\mathbb R$, nondecreasing with
$\Phi_{\mathrm{HI}}^{\widehat T,g(0)}(s)/s$ decreasing to zero for $s>0$,
such that
\[
 \operatorname{Rm}_{g(t)}\ge
 -\Phi_{\mathrm{HI}}^{\widehat T,g(0)}(R_{g(t)})
 \qquad(0\le t\le\widehat T).
\]
The producer is allowed to depend on the horizon and on the time-zero metric;
every later use fixes those data before invoking the estimate.
\end{lemma}

\begin{proof}
The Hamilton--Ivey inequality in
\ref{def:hamilton-ivey-pinching-condition}, applied on the compact interval
$[0,\widehat T]$, gives a bound depending on $\widehat T$ and the lower
curvature bound of $g(0)$.  The construction following Kleiner--Lott
Appendix~B.5--B.6 smooths that bound to a nondecreasing function with the
required sublinear quotient and enlarges it on a compact range to include
$t=0$.  The maximum principle then gives the displayed inequality for every
$t$ in this fixed interval.  No assertion that one producer works for all
finite horizons is made.
\end{proof}

\begin{lemma}[Imported interface: pinching-preserving cap insertion]
\label{lem:ricci-flow-surgery-cap-insertion}
\dcref{KL:II.4:72.24}
\uses{lem:standard-surgery-cap-model,def:standard-solution,
  def:hamilton-ivey-pinching-condition,def:epsilon-neck-and-scale}
\source{kleiner-lott}
\notready
For the fixed $A=10$ and $(B,F,k,g_0)$ above, there are
$\delta_0^{\rm cap}\in(0,1/2]$ and $\delta'(\delta)\to0$ such that, whenever
$t>0$, $0<\delta<\delta_0^{\rm cap}$, and a metric $h_0$ on
$(-A,\delta^{-1})\times S^2$ is $\delta$-close after scaling by $R(x)$ to the
unit scalar-curvature cylinder in $C^{\lfloor1/\delta\rfloor+1}$, with a
specified base point $x\in\{0\}\times S^2$, and satisfies the time-$t$
Hamilton--Ivey condition on the whole input metric, it extends over a
three-ball to a metric $h$ with the same pinching,
$h=h_0$ on $[0,\delta^{-1})\times S^2$, $R(x)h=g_0$ on
$(-B,-A)\times S^2$, constant curvature $k^2$ on the ball, and
$\delta'$-closeness to $e^{2F}g_{\mathrm{cyl}}$ on the capped end.
\end{lemma}

\begin{proof}
Rescale the input neck by $R(x)^{1/2}$ and write it as a
$C^{\lfloor1/\delta\rfloor+1}$-small perturbation of the unit cylinder on
$(-A,\delta^{-1})\times S^2$.  Insert the fixed profile from
\ref{lem:standard-surgery-cap-model} on the left, using a cutoff only where
the displayed bounds on $f,f',f''$ hold.  The warped-product curvature
formulas bound every cutoff error by $c\delta$; choosing
$\delta_0^{\rm cap}$ small makes the Hamilton--Ivey cone invariant under
these errors.  The cap is the rescaled standard solution on $(-B,-A)$ and is
constant-curvature on the closing ball, while the right-hand cylinder is
unchanged.  Scaling back gives the factor $R(x)$ and a residual neck error
$\delta'(\delta)\to0$.  This is the pinching-preserving cap insertion of
Kleiner--Lott II.4,72.24.
\end{proof}

\begin{lemma}[Finite-horizon Hamilton--Ivey envelope across surgery]
\label{lem:finite-horizon-surgery-phi-envelope}
\dcref{KL:appB:B.5,KL:appB:after-B.6,KL:II.4:72.24}
\source{kleiner-lott}
\uses{def:ricci-flow-with-surgery,lem:hamilton-ivey-phi-envelope,
  lem:ricci-flow-surgery-cap-insertion}
\notready
Let $(\mathcal M,G)$ be a three-dimensional Ricci flow with surgery on a
finite interval $[0,b]$ with finitely many transition times.  Assume that
the initial curvature operator has smallest eigenvalue at least $-1$, every
ordinary piece has bounded curvature on compact time intervals, and every
transition is made by the pinching-preserving cap insertion of
\ref{lem:ricci-flow-surgery-cap-insertion}.  Then there is one admissible
smooth nondecreasing function $\Phi_{\rm HI}^{b,g(0)}$ with
$\Phi_{\rm HI}^{b,g(0)}(s)/s\to0$ as $s\to\infty$ such that
\[
 \operatorname{Rm}_{G}(x,t)\ge-\Phi_{\rm HI}^{b,g(0)}(R_G(x,t))
\]
on every ordinary slice and on both the backward and forward sides of every
transition where those slices are defined.
\end{lemma}

\begin{proof}
Apply the finite-horizon ordinary-flow envelope
\ref{lem:hamilton-ivey-phi-envelope} on each of the finitely many smooth
pieces, using the incoming lower-curvature bound supplied by the preceding
cap insertion.  At a transition, the cap lemma gives the same Hamilton--Ivey
inequality on the inserted metric, so no new lower-curvature branch is
created.  Take the pointwise maximum of the finitely many envelopes and
smooth it upward on a compact range; the maximum remains nondecreasing and
sublinear, and the resulting single function controls every piece.  This is
the finite-induction construction in Kleiner--Lott Appendix~B and II.4,
not an assertion that the ordinary-flow producer is uniform over arbitrary
horizons.
\end{proof}

\begin{definition}[Kappa-solution]
\label{def:kappa-solution}
\dcref{MT:ch9:9.1,MT:ch9:9.2,KL:I.11:38.1}
\source{morgan-tian}
\notready
For $d\in\{2,3\}$, a $d$-dimensional \emph{$\kappa$-solution} is a connected Ricci flow
$(M,g(t))$, $-\infty<t\le0$, such that every time slice is complete and
nonflat, has bounded nonnegative curvature operator, and is
$\kappa$-noncollapsed on every finite scale.  The value of $\kappa>0$ is
fixed for the flow.
\end{definition}

\begin{definition}[Kleiner--Lott canonical-neighborhood assumption]
\label{def:ricci-flow-surgery-canonical-neighborhood-assumption}
\dcref{KL:I.11:50.1,KL:II.1:59.4,KL:II.1:59.7,KL:II.2:Claim5,KL:II.4:69.1}
\source{kleiner-lott}
\uses{def:ricci-flow-with-surgery,def:kappa-solution,def:standard-solution,
  def:epsilon-neck-and-scale,def:canonical-cap}
\notready
Fix small $\epsilon>0$, the universal $\kappa_0$ of Proposition~50.1, the
derivative constant $\eta$ of Lemma~59.4, and the constants
$C_1^\kappa,C_2^\kappa,C_1^{\mathrm{std}},C_2^{\mathrm{std}}$ from
Lemma~59.7 and Claim~5 of II.2.  Set
\[
 C_1=\max\{30\epsilon^{-1},C_1^\kappa,C_1^{\mathrm{std}}\},\qquad
C_2=\max\{C_2^\kappa,C_2^{\mathrm{std}}\}.
\]
At a transition time, $\mathcal M_t^-$ denotes the presurgery slice and
$\mathcal M_t^+$ the postsurgery slice; at a regular time both denote the
ordinary slice.  A backward (respectively forward) parabolic neighborhood
below is only formed with $x\in\mathcal M_t^-$ (respectively
$x\in\mathcal M_t^+$) and with time depth at most $t$.  With this eligibility
convention, for a positive nonincreasing $r$, every
$(x,t)\in\mathcal M_t^\pm$ with
$R(x,t)\ge r(t)^{-2}$ has an open canonical neighborhood
$\overline{B^\pm(x,t,\widehat r)}\subset U\subset B^\pm(x,t,2\widehat r)$,
where $R(x,t)^{-1/2}<\widehat r<C_1R(x,t)^{-1/2}$, modeled by a strong
$\epsilon$-neck, a $\kappa_0$- or standard-solution cap, an $S^3/\mathbb{RP}^3$
component, or an $\epsilon$-round positive space form; on $U$ scalar curvature
is $C_2$-comparable to $R(x,t)$ and the derivative bounds
$|\nabla R|<\eta R^{3/2}$ and $|\partial_tR|<\eta R^2$ hold one-sided at surgery.
In the neck, cap, and closed-component cases,
\[
 \operatorname{Vol}_{g(t)}U>C_2^{-1}R(x,t)^{-3/2},
\]
and in the closed-component case the infimal sectional curvature on $U$ is
greater than $C_2^{-1}R(x,t)$.
\end{definition}

\begin{definition}[Kleiner--Lott $\Phi$-pinching]
\label{def:ricci-flow-surgery-phi-pinching}
\dcref{KL:II.4:69.5}
\source{kleiner-lott}
\uses{def:ricci-flow-with-surgery}
\notready
For an admissible smooth positive nondecreasing $\Phi$ with $\Phi(s)/s$ decreasing
to zero on $(0,\infty)$, the $\Phi$-pinching assumption is
$\operatorname{Rm}(x,t)\ge-\Phi(R(x,t))$ at every spacetime point.
\end{definition}

\begin{definition}[Kleiner--Lott a-priori assumptions]
\label{def:ricci-flow-surgery-a-priori-assumptions}
\dcref{KL:II.4:69.6}
\source{kleiner-lott}
\uses{def:ricci-flow-surgery-canonical-neighborhood-assumption,
  def:ricci-flow-surgery-phi-pinching}
\notready
A flow on $I\subset[0,\infty)$ satisfies the a-priori assumptions for $(r,\Phi)$
when it satisfies both the preceding $r$-canonical-neighborhood and
$\Phi$-pinching assumptions on $I$.
\end{definition}

\begin{lemma}[Imported interface: Kleiner--Lott surgery scale]
\label{lem:ricci-flow-surgery-horn-scale}
\dcref{KL:II.4:71.1,KL:II.4:after-71.1}
\uses{def:ricci-flow-surgery-a-priori-assumptions,def:epsilon-horn-terminology,
  def:embedded-horn-frontier,def:canonical-neck-cap-regions}
\source{kleiner-lott}
\notready
Fix $\Phi$, $\widehat T$, and positive nonincreasing $r$ on
$[0,\widehat T]$.  At a singular time $T$, write $\bar R$ for the scalar
curvature on the terminal maximal extension and $\bar R(x,T)$ for its value
at a point of the bounded-curvature domain.  There is a producer
\[
 H(\Phi,T,r,\delta)>0\qquad(0<T\le\widehat T,\ 0<\delta<1/2)
\]
with $H(\Phi,T,r,\delta)<\delta^2r(T)$, nondecreasing in $r(T)$ and $\delta$,
such that at a singular time $T$, an embedded $\epsilon$-horn whose compact
bounded-curvature frontier lies in $\{\bar R\le(\delta r(T))^{-2}\}$ and whose
point has $\bar R\ge H(\Phi,T,r,\delta)^{-2}$ contains the backward parabolic
region
\[
 P\bigl(x,T,\delta^{-1}\bar R(x,T)^{-1/2},-\bar R(x,T)^{-1}\bigr)
\]
inside a strong $\delta$-neck.
Moreover, for every $r_->0$ and $\delta_->0$ there is a uniform lower
envelope
\[
 H_*=H_*(\Phi,\widehat T,r_-,\delta_-)>0
\]
such that $H(\Phi,T,r,\delta)\ge H_*$ whenever
$0<T\le\widehat T$, $r(T)\ge r_-$, and $\delta\ge\delta_-$.  Thus a
positive nonincreasing control function on a bounded interval cannot make the
selected surgery scales collapse to zero.
\end{lemma}

\begin{proof}
The horn-scale producer is the Kleiner--Lott scale-selection theorem.  On a
bounded stage its constants are chosen uniformly for $T\in[0,\widehat T]$.
The monotonicity in the lower canonical-neighborhood parameter and in the
control parameter therefore gives the displayed envelope after substituting
$r_-$ and $\delta_-$.  Since a positive nonincreasing $\delta$ satisfies
$\delta(t)\ge\delta(\widehat T)$ on this interval, the envelope is positive
for every surgery time.  This is the source's bounded-stage scale estimate,
not an additional regularity assumption on $H$.
\end{proof}

\begin{definition}[Ricci flow with Kleiner--Lott $(r,\delta)$-cutoff]
\label{def:ricci-flow-with-cutoff}
\dcref{KL:II.4:73.1}
\source{kleiner-lott}
\uses{def:ricci-flow-with-surgery,def:ricci-flow-surgery-a-priori-assumptions,
  lem:finite-horizon-surgery-phi-envelope,lem:ricci-flow-surgery-horn-scale,
  lem:ricci-flow-surgery-cap-insertion,def:epsilon-horn-terminology}
\notready
On $[0,b]$, let $r$ be positive nonincreasing, let
$\bar\delta:[0,b]\to(0,\infty)$ be nonincreasing, and fix a positive
nonincreasing threshold sequence
$\Delta=(\delta_i)_{i\ge0}$ with $0<\delta_i\le\delta_0^{\rm cap}$, where
$\delta_0^{\rm cap}$ is the cap-insertion threshold supplied by
\ref{lem:ricci-flow-surgery-cap-insertion}.  Put
$t_0=2^{-5}$ and $T_i=2^it_0$.  Require that there is no transition on the
initial short-time interval $[0,T_0]$ and that every later transition obeys
\[
\bar\delta(t)<\delta_{i+1}\le\delta_0^{\rm cap}
 \qquad(T_i<t\le T_{i+1}).
\]
Fix the envelope $\Phi_{\mathrm{HI}}^{b,g(0)}$ supplied by
\ref{lem:finite-horizon-surgery-phi-envelope}.  At each transition set
$h(t)=H(\Phi_{\mathrm{HI}}^{b,g(0)},t,r,\bar\delta(t))$ and put
$\rho=\bar\delta(t)r(t)$.  Identify the final smooth interval before this
transition with $M_{T^-}$, where $T=t$, and define
\[
 \Omega=\{x\in M_{T^-}:\limsup_{s\nearrow T}|\operatorname{Rm}(x,s)|<\infty\}.
\]
The limiting metric on $\Omega$ has scalar curvature $\bar R$, the function
$\bar R$ is proper and bounded below, and
$\Omega_\rho=\{y\in\Omega:\bar R(y)\le\rho^{-2}\}$ is compact.  Retain
exactly the connected components of $\Omega$ meeting $\Omega_\rho$ and discard
the other components.  For every remaining embedded $\epsilon$-horn $\mathcal H$
whose compact bounded-curvature frontier lies in $\Omega_\rho$, choose (using
the properness of $\bar R$) a point $x_{\mathcal H}$ on the horn with
$\bar R(x_{\mathcal H})=h(t)^{-2}$.  The scale lemma gives a strong
$\bar\delta(t)$-neck about this point; cut along its central cross-sectional
$2$-sphere and glue the cap of
\ref{lem:ricci-flow-surgery-cap-insertion} along that sphere.  Thus the cut
is made at a scale-selected neck, not at an arbitrary frontier of
$\Omega_\rho$.  The whole backward slice satisfies the compatible time-$t$
Hamilton--Ivey condition.  On every open smooth subinterval ending at
a transition, the flow also satisfies the a-priori $(r,\Phi_{\mathrm{HI}})$
conditions of
\ref{def:ricci-flow-surgery-a-priori-assumptions}; at a first failure the
canonical-neighborhood condition is required only on the half-open interval
before the failure time.  An all-time flow has this cutoff when every finite
restriction has the stated property.  The condition is imposed only where
surgery occurs, so the displayed stagewise bounds are vacuous on a
surgery-free stage.  This is a contract for an already constructed flow;
the extension theorem below produces parameters making these a-priori
conditions hold.
\end{definition}

\begin{definition}[Strong canonical-neighborhood control]
\label{def:strong-canonical-neighborhood-control}
\dcref{MT:ch9:9.76,MT:ch14:14.18}
\source{morgan-tian}
\uses{def:ricci-flow-with-surgery,def:epsilon-neck-and-scale,def:canonical-cap}
\notready
Fix $C<\infty$, $0<\epsilon<1/2$, and $r>0$.  A three-dimensional Ricci
flow with surgery satisfies the \emph{strong
$(C,\epsilon)$-canonical-neighborhood assumption with parameter $r$} if
every point $x$ with $R(x)\ge r^{-2}$ has one of the following
neighborhoods at the curvature scale $R(x)^{-1/2}$:
\begin{enumerate}
\item a strong $\epsilon$-neck centered at $x$, whose parabolic rescaling by
      $R(x)$ is $\epsilon$-close in
      $C^{\lfloor\epsilon^{-1}\rfloor}$ to the shrinking round cylinder on
      normalized times $(-1,0]$;
\item a $(C,\epsilon)$-cap whose core contains $x$, with $C$-comparable
      scalar curvatures and scale-normalized diameter and volume bounded by
      $C$, together with a uniform core volume lower bound and the
      scale-invariant derivative bounds for $R$;
\item a compact positively curved $C$-component containing $x$, with
      underlying manifold diffeomorphic to $S^3$ or $\mathbb{RP}^3$,
      least sectional curvature greater than $C^{-1}\sup_M R$, and
      \[
       C^{-1}(\inf_M R)^{-1/2}<\operatorname{diam}(M)
       <C(\sup_M R)^{-1/2};
      \]
\item an $\epsilon$-round component containing $x$: a compact connected
      component $(M,g)$ for which there are $R>0$, a compact
      constant-curvature-one manifold $(Z,g_0)$, and a diffeomorphism
      $\varphi:Z\to M$ such that $\varphi^*(Rg)$ is $\epsilon$-close to $g_0$
      in the $C^{\lfloor\epsilon^{-1}\rfloor}$ topology.
\end{enumerate}
Only the first alternative is a spacetime neighborhood.  The other
alternatives lie in the slice $M_{\mathbf t(x)}$.
\end{definition}

\begin{definition}[Curvature pinched toward positive]
\label{def:curvature-pinched-toward-positive}
\dcref{MT:ch10:10.1,MT:ch14:14.17}
\source{morgan-tian}
\uses{def:ricci-flow-with-surgery}
\notready
Let $(\mathcal M,G)$ be a three-dimensional Ricci flow with surgery whose
time domain lies in $[0,\infty)$.  At $x\in\mathcal M$, write
$\lambda(x)\ge\mu(x)\ge\nu(x)$ for the eigenvalues of the curvature
operator on $\bigwedge^2\mathcal H_x$ and set
$X(x)=\max\{0,-\nu(x)\}$.  Its curvature is \emph{pinched toward positive}
when, for every $x$,
\[
 R(x)\ge-\frac{6}{1+4\mathbf t(x)}
\]
and, whenever $X(x)>0$,
\[
 R(x)\ge
 2X(x)\bigl(\log X(x)+\log(1+\mathbf t(x))-3\bigr).
\]
For a metric at a fixed time $T\ge0$, pinching up to time $T$ means the
same two inequalities with $\mathbf t(x)$ replaced by $T$.
\end{definition}

\begin{proposition}[Imported interface: gluing evolving necks]
\label{prop:evolving-neck-gluing}
\dcref{MT:ch15:15.1}
\uses{def:epsilon-neck-and-scale}
\source{morgan-tian}
\notready
There is $0<\beta<1/2$ such that the following holds for every
$0<\epsilon<1$.  Let
$(N\times[-t_0,0],g_1(t))$ be an evolving
$\beta\epsilon$-neck centered at $x$ with $R(x,0)=1$, and let
$(N'\times(-t_1,-t_0],g_2(t))$ be a strong
$\beta\epsilon/2$-neck.  Suppose an isometric identification of their
time-$-t_0$ slices sends $(x,-t_0)$ to the center of the strong neck.
Then the induced flow on
\[
 N\times[-t_0,0]\mathbin{\cup}N'\times(-t_1,-t_0]
\]
contains a strong $\epsilon$-neck centered at $(x,0)$.
\end{proposition}

\begin{proof}
Rescale the first neck by its centre curvature.  On the overlap slice at
$-t_0$, the two parametrizations are $C^{[1/\epsilon]}$-close; after a
time-independent cylindrical reparametrization their difference is
$C^{[1/\epsilon]-1}$-small.  The curvature and derivative estimates on the
two parabolic regions propagate this error to time $0$ by the parabolic
maximum principle.  Strong-neck stability then extends the first neck across
the overlap, and the choice $\beta<1/2$ makes the total error less than
$\epsilon$.  Undoing the rescaling gives a strong $\epsilon$-neck on the
displayed union.  This is the overlap proof of Morgan--Tian Proposition 15.1.
\end{proof}

\begin{definition}[Controlled surgery parameters]
\label{def:controlled-surgery-parameters}
\dcref{MT:ch15:15.5,MT:ch15:15.6,MT:ch15:15.7}
\source{morgan-tian}
\uses{def:strong-canonical-neighborhood-control,prop:evolving-neck-gluing,
  thm:local-neck-surgery,lem:ricci-flow-surgery-horn-scale,
  lem:finite-horizon-surgery-phi-envelope,lem:ricci-flow-surgery-cap-insertion}
\notready
On every finite restriction $[0,b]$ of the all-time construction, fix the
finite-horizon surgery producer $\Phi_{\mathrm{HI}}^{b,g(0)}$ from
\ref{lem:finite-horizon-surgery-phi-envelope}; all occurrences of
$\Phi_{\mathrm{HI}}$ below mean this fixed producer on that restriction.
Fix $\epsilon>0$ sufficiently small and choose $C$ larger than the uniform
canonical-neighborhood constants for the standard solution and for
three-dimensional $\kappa$-solutions.  Put $t_0=2^{-5}$ and
$T_i=2^it_0$.  The control data are positive nonincreasing sequences
\[
 \mathbf r=(r_i)_{i\ge0},\qquad
 \mathbf K=(\kappa_i)_{i\ge0},\qquad
 \boldsymbol\delta=(\delta_i)_{i\ge0},
\qquad r_0=\epsilon,
\]
where $\kappa_0>0$ and $0<\delta_0\le\delta_0^{\rm cap}$ are fixed initial
entries.  The all-time
existence theorem chooses $\kappa_0$ from the normalized short-time estimate
and chooses $\delta_0$ below the fixed local-surgery, neck-gluing, and
cap-barrier thresholds.  Their associated step functions are
\[
 r(t)=r_0,\quad\kappa(t)=\kappa_0
 \quad(0\le t<T_0),
\]
and have values $r_{i+1}$ and $\kappa_{i+1}$ on
$[T_i,T_{i+1})$.  A positive nonincreasing surgery-control function
$\bar\delta(t)$ is required to have no transition on $[0,T_0]$ and to
satisfy, at every later transition time,
\[
 \bar\delta(t)<\delta_{i+1}\le\delta_0^{\rm cap}
 \qquad(T_i<t\le T_{i+1}),
\]
where $\delta_0^{\rm cap}$ is the cap-insertion threshold in
\ref{lem:ricci-flow-surgery-cap-insertion}.
Define
\[
 \rho(t)=\bar\delta(t)r(t),
 \qquad
 h(t)=H(\Phi_{\mathrm{HI}}^{b,g(0)},t,r,\bar\delta(t))
 \le \bar\delta(t)^2r(t),
\]
at surgery times; this is the neck scale supplied by
\ref{lem:ricci-flow-surgery-horn-scale}.
At a surgery time $t$, every cutting sphere is the central sphere of a strong
$\bar\delta(t)$-neck of scalar curvature $h(t)^{-2}$.
\end{definition}

\begin{definition}[Ricci flow with controlled surgery data]
\label{def:controlled-ricci-flow-with-surgery}
\dcref{MT:ch15:15.2,MT:ch15:15.5,MT:ch15:15.8}
\source{morgan-tian}
\uses{def:ricci-flow-with-surgery,def:strong-canonical-neighborhood-control,def:controlled-surgery-parameters,def:maximal-ricci-flow-extension}
\notready
For the fixed $C$, $\epsilon$, and control functions, a
\emph{controlled Ricci flow with surgery} is a three-dimensional Ricci flow
with surgery satisfying Assumptions~(1)--(7) of Morgan--Tian and using the
prescribed cutting data.  Explicitly:
\begin{enumerate}
\item Every slice is compact and contains no embedded $\mathbb{RP}^2$ with
      trivial normal bundle.
\item Singular times form a discrete subset of the time interval.
\item Time zero is initial, and every component of $(M_0,g(0))$ is
      normalized: $|\operatorname{Rm}|\le1$ and every unit ball has at least
      half the Euclidean unit-ball volume.
\item At a singular time $t$, the exposed region $E_t$ is a finite disjoint
      union of smoothly embedded closed three-balls.  Its complement
      $C_t=M_t\setminus\operatorname{int}E_t$ is the maximal continuing
      region: it is the maximal subset for which the time vector field gives
      an embedding $C_t\times(t-\eta,t]\to\mathcal M$ for some $\eta>0$.
      Thus
      $M_t=C_t\cup_{\Sigma_t}E_t$ with
      $\Sigma_t=\partial C_t=\partial E_t$ a finite union of two-spheres.
      For every sufficiently close regular $t^-<t$, backward flow maps $C_t$
      diffeomorphically onto a compact codimension-zero submanifold
      $C_{t^-}\subset M_{t^-}$ with boundary the backward image of
      $\Sigma_t$; the disappearing region is
      $D_{t^-}=M_{t^-}\setminus C_{t^-}$.
\item Every component of $\Sigma_t$ is the central sphere of a strong
      $\bar\delta(t)$-neck at scalar curvature $h(t)^{-2}$ in the maximal
      presurgery extension to $t$.
\item For every singular time $t$ and every sufficiently close regular
      $t^-<t$, every point of the disappearing region
      $D_{t^-}=M_{t^-}\setminus C_{t^-}$ has a strong
      $(C,\epsilon)$-canonical neighborhood.
\item Between consecutive singular times the ordinary Ricci flow is maximal.
\end{enumerate}
Positive pinching, the strong canonical-neighborhood assumption with parameter
$r(t)$, and $\mathbf K$-noncollapsing are separate quantitative properties;
they are not part of this definition.
\end{definition}

\begin{theorem}[Local volume and injectivity-radius control]
\label{thm:local-volume-injectivity-radius-control}
\dcref{MT:ch1:1.35,MT:ch1:1.36,MT:ch5:5.10}
\source{morgan-tian}
\notready
Fix $n$, $K<\infty$, $r_0>0$, and $v_0>0$.  There are positive constants
\[
 v_*=v_*(n,K,r_0,v_0),
 \qquad \iota_*=\iota_*(n,K,r_0,v_0),
\]
with the following property.  If $(M^n,g)$ is complete,
$|\operatorname{Rm}|\le K$ on $B(p,r_0)$, and
$\operatorname{Vol}B(p,r_0)\ge v_0$, then, for
$0<s\le\min\{r_0,K^{-1/2}\}$,
\[
 \operatorname{Vol}B(p,s)\ge v_*s^n,
 \qquad
 \operatorname{inj}_g(p)\ge
 \iota_*\min\{r_0,K^{-1/2}\}.
\]
The convention is that $K^{-1/2}=\infty$ when $K=0$.
\end{theorem}

\begin{proof}
Put $a=\min\{r_0,K^{-1/2}\}$ and rescale by $a^{-2}$, so the relevant
ball has radius one and sectional curvatures have absolute value at most one.
In particular $\operatorname{Ric}\ge-(n-1)$.  Bishop--Gromov comparison with
hyperbolic $n$-space gives, for $0<s\le1$,
\[
 \frac{\operatorname{Vol}B(p,s)}{V_{-1}(s)}
 \ge
 \frac{\operatorname{Vol}B(p,1)}{V_{-1}(1)}.
\]
Applying the same comparison first between radii $a$ and $r_0$ before
rescaling supplies a positive lower bound for
$\operatorname{Vol}B(p,1)$ depending only on the displayed data.  Since
$V_{-1}(s)/s^n$ has a positive minimum on $(0,1]$, this proves the asserted
volume ratio.

For the injectivity estimate, curvature at most one excludes conjugate points
on radial geodesics of length less than one.  If the injectivity radius at
$p$ tended to zero with the volume ratio fixed, Klingenberg's alternative
would therefore give a shortest geodesic loop of length twice the injectivity
radius.  Lift the unit ball to the nonsingular exponential ball in $T_pM$.
The successive lifts obtained by concatenating this loop are distinct until
their total length reaches a fixed fraction of one.  Curvature comparison
bounds from below the volume of one fixed lifted subball, while the distinct
lifts all project into $B(p,1)$ with uniformly bounded multiplicity.  Hence
the volume of $B(p,1)$ is at most a dimensional constant times the
injectivity radius, contradicting the positive volume lower bound.  This
gives a positive rescaled injectivity bound; scaling back gives $\iota_*a$.
\end{proof}

\begin{theorem}[Normalized short-time flow]
\label{thm:normalized-short-time-ricci-flow}
\dcref{MT:ch3:3.11,MT:ch4:4.11,MT:ch5:5.16,MT:ch15:15.1}
\source{morgan-tian}
\notready
There is a universal $\kappa_0>0$ such that every normalized compact
three-manifold $(M,g_0)$ admits a unique Ricci flow $g(t)$ on
$[0,2^{-5}]$ with $g(0)=g_0$.  On this interval
\[
 |\operatorname{Rm}_{g(t)}|\le2,
 \qquad
 \operatorname{Vol}_{g(t)}B_{g(t)}(x,\rho)\ge\kappa_0\rho^3
 \quad(0<\rho\le1).
\]
Consequently this interval has no surgery time and supplies the initial
stage of a controlled Ricci flow with surgery.
\end{theorem}

\begin{proof}
\uses{def:ricci-flow-with-surgery,def:controlled-ricci-flow-with-surgery,
  def:maximal-ricci-flow-extension,def:normalized-riemannian-three-manifold,
  thm:local-volume-injectivity-radius-control}
The Ricci--DeTurck equation obtained by adding the Lie derivative along the
trace of the difference between the connections of $g$ and $g_0$ is strictly
parabolic.  Short-time parabolic existence, followed by pullback under the
generating diffeomorphisms, gives a Ricci flow; the reverse gauge argument
gives uniqueness.

The curvature evolution inequality and $|\operatorname{Rm}_{g_0}|\le1$ give
$|\operatorname{Rm}_{g(t)}|\le2$ for $t\le2^{-4}$ as long as the flow
exists.  On the same interval metric and volume distortion give universal
constants $a,b>0$ such that
\[
 B_{g_0}(x,a\rho)\subset B_{g(t)}(x,\rho),
 \qquad
 d\operatorname{vol}_{g(t)}\ge b\,d\operatorname{vol}_{g_0}.
\]
Apply \ref{thm:local-volume-injectivity-radius-control} at time zero.
Normalization and $|\operatorname{Rm}_{g_0}|\le1$ give one universal lower
volume ratio $v_*>0$ on all radii at most one.  The distortion estimates therefore
yield
\[
 \operatorname{Vol}_{g(t)}B_{g(t)}(x,\rho)
 \ge b\operatorname{Vol}_{g_0}B_{g_0}(x,a\rho)
 \ge bv_*a^3\rho^3.
\]
Take $\kappa_0=bv_*a^3$.  If the maximal flow ended before
$2^{-5}$, its uniform curvature bound and the bounded-curvature extension
theorem would extend it, a contradiction.  Thus it exists through $2^{-5}$;
viewed as a surgery flow, its surgery set is empty.
\end{proof}

\section{Reduced Geometry}

\begin{theorem}[Local Ricci-flow derivative estimates]
\label{thm:local-ricci-flow-derivative-estimates}
\dcref{MT:ch3:3.28,MT:ch3:3.33}
\source{morgan-tian}
\notready
For every $m\ge0$ there is $C_m<\infty$ with the following property.  If a
Ricci flow has a relatively compact parabolic neighborhood
$P(x,t,r,-r^2)$ on which $|\operatorname{Rm}|\le r^{-2}$, then
\[
 |\nabla^m\operatorname{Rm}|
 \le C_m r^{-m-2}
\]
on $P(x,t,r/2,-r^2/2)$.
\end{theorem}

\begin{proof}
\uses{def:ricci-flow-with-surgery}
Parabolically rescale so that $r=1$.  Commuting covariant derivatives with
the curvature evolution equation gives
\[
 (\partial_t-\Delta)\nabla^m\operatorname{Rm}
 =\sum_{a+b=m}\nabla^a\operatorname{Rm}*\nabla^b\operatorname{Rm}.
\]
For $m=1$, apply the maximum principle to
$t|\nabla\operatorname{Rm}|^2+A|\operatorname{Rm}|^2$ after multiplying by
a spatial cutoff supported in the unit ball; choosing $A$ larger than the
universal reaction coefficients absorbs the cross terms.  This gives a
uniform first-derivative bound on the half cylinder.

Inductively, assume the estimates through order $m-1$.  Apply the same
argument to
\[
 t^m|\nabla^m\operatorname{Rm}|^2
 +\sum_{j<m}A_jt^j|\nabla^j\operatorname{Rm}|^2,
\]
using nested cutoffs.  The induction bounds every lower-order product, and
the constants $A_j$ absorb the commutator terms.  The parabolic maximum
principle gives a bound depending only on $m$.  Scaling back yields the
factor $r^{-m-2}$.
\end{proof}

\begin{definition}[Reduced length]
\label{def:reduced-length}
\dcref{MT:ch6:6.2,MT:ch6:6.45}
\source{morgan-tian}
\uses{def:ricci-flow-with-surgery}
\notready
Fix a smooth point $(x,T)$ and write backward time as $\tau=T-t$.  For a
piecewise smooth backward path $\gamma\colon[0,\bar\tau]\to\mathcal M$
with $\gamma(0)=x$ and $\mathbf t(\gamma(\tau))=T-\tau$, put
$X_\gamma=\dot\gamma+\chi\in\mathcal H$.  Define
\[
 \mathcal L(\gamma)=
 \int_0^{\bar\tau}\sqrt\tau
 \bigl(R(\gamma(\tau))+|X_\gamma(\tau)|_G^2\bigr)\,d\tau.
\]
For a point $y$ at backward time $\bar\tau$, let
$L_x(y,\bar\tau)$ be the infimum of $\mathcal L$ over such paths and put
\[
 l_x(y,\bar\tau)=\frac{L_x(y,\bar\tau)}{2\sqrt{\bar\tau}}.
\]
These quantities are used only on domains on which every competing path
stays in the smooth locus.
\end{definition}

For estimates in which a surgery cap is treated as a barrier, we also use the
positive-part version of reduced length.  Along a path in the same domain set
\[
 R_+(z):=\max\{R(z),0\},\qquad
 \mathcal L_+(\gamma):=
 \int_0^{\bar\tau}\sqrt\tau\bigl(R_+(\gamma(\tau))+
 |X_\gamma(\tau)|_G^2\bigr)\,d\tau.
\]
Thus $\mathcal L\le\mathcal L_+$ pointwise.  When
$R\ge-6$ on the domain and $\bar\tau\le\widehat\tau$, the elementary
comparison
\[
 0\le \mathcal L_+(\gamma)-\mathcal L(\gamma)
 \le 4\widehat\tau^{3/2}
\]
is the bounded-error interface used below (Morgan--Tian, Chapters~6 and~15).

\begin{definition}[Reduced exponential map and injectivity domain]
\label{def:reduced-exponential-map}
\dcref{MT:ch6:6.17,MT:ch6:6.25,MT:ch6:6.26}
\source{morgan-tian}
\uses{def:reduced-length}
\notready
An $\mathcal L$-geodesic from $(x,T)$ is a critical point of
$\mathcal L$.  Its initial vector is
\[
 v=\lim_{\tau\downarrow0}\sqrt\tau\,X_\gamma(\tau)
 \in T_xM_T.
\]
The reduced exponential map is
$\mathcal L\exp_x^\tau(v)=\gamma_v(\tau)$ whenever this geodesic exists.
Its injectivity set $\mathcal D_x^\tau$ consists of those $v$ for which
$\gamma_v|_{[0,\tau]}$ is the unique minimizing $\mathcal L$-geodesic to its
endpoint and has no nonzero $\mathcal L$-Jacobi field vanishing at both ends.
The image $\Omega_x^\tau=\mathcal L\exp_x^\tau(\mathcal D_x^\tau)$ is the
reduced injectivity domain.  To match Morgan--Tian's notation in the
stagewise statements below, write
$\widetilde{\mathcal U}_x(\tau):=\mathcal D_x^\tau$.
\end{definition}

\begin{proposition}[Monotonicity of the reduced injectivity sets]
\label{prop:reduced-injectivity-set-monotonicity}
\dcref{MT:ch6:6.30}
\source{morgan-tian}
\uses{def:reduced-exponential-map}
\notready
If $0<\tau<\tau'$, then
\[
 \mathcal D_x^{\tau'}\subset\mathcal D_x^\tau.
\]
Thus an initial-vector set contained in $\mathcal D_x^{\tau_0}$ has
well-defined injective prefixes at every earlier time.
\end{proposition}

\begin{proof}
For $v\in\mathcal D_x^{\tau'}$, the minimizing $\mathcal L$-geodesic
to $\mathcal L\exp_x^{\tau'}(v)$ has minimizing prefixes, so its restriction
to $[0,\tau]$ is minimizing.  If the differential of
$\mathcal L\exp_x^\tau$ were singular, a nonzero $\mathcal L$-Jacobi field
would vanish at both endpoints.  Taking the first singular time and extending
that Jacobi field by zero along the remaining tail gives a broken variation;
the second-variation characterization of an $\mathcal L$-minimizer would
then force this nonsmooth extension to be Jacobi, a contradiction.  If local
uniqueness failed at the prefix, concatenating a second minimizing prefix
with the original tail would give a nonsmooth path no longer than the
$\tau'$-minimizer, again a contradiction.  Hence both injectivity conditions
persist to every earlier time.
\end{proof}

\begin{definition}[Reduced volume]
\label{def:reduced-volume}
\dcref{MT:ch6:6.70,MT:ch6:6.71}
\source{morgan-tian}
\uses{def:reduced-length,def:reduced-exponential-map}
\notready
For a measurable set $A$ in the regular reduced-exponential image at time
$T-\tau$, its reduced volume based at $(x,T)$ is
\[
 \widetilde V_x(A,\tau)=
 \int_A(4\pi\tau)^{-3/2}e^{-l_x(y,\tau)}
 \,d\operatorname{vol}_{g(T-\tau)}(y).
\]
Equivalently, if $J(v,\tau)$ is the Jacobian of the reduced exponential map
and $\widetilde A=(\mathcal L\exp_x^\tau)^{-1}(A)$, then
\[
 \widetilde V_x(A,\tau)=
 \int_{\widetilde A}(4\pi\tau)^{-3/2}
 e^{-l_x(\mathcal L\exp_x^\tau(v),\tau)}J(v,\tau)\,dv.
\]
When the regular image has full measure in the slice, write
$\widetilde V_x(\tau)$ for the value obtained by taking that whole image.
For a family $A(\tau)=\mathcal L\exp_x^\tau(\widetilde A)$ with
$0<\tau\le\tau_0$, the reduced-volume monotonicity assertion below is used
only when there is one open smooth spacetime set $\mathcal O$ containing
$x$ and every $A(\tau)$ such that every minimizing $\mathcal L$-geodesic
from $x$ to a point of $A(\tau)$ stays in $\mathcal O\cup\{x\}$.
This is the common-domain hypothesis in Morgan--Tian's monotonicity theorem.
\end{definition}

\begin{lemma}[Upper support inequality for reduced length]
\label{lem:reduced-length-support-inequality}
\dcref{MT:ch6:6.4,MT:ch6:6.13,MT:ch6:6.49,MT:ch6:6.51,MT:ch7:7.8}
\source{morgan-tian}
\notready
On a smooth bounded-curvature Ricci-flow domain, reduced length is smooth
before the reduced cut locus and satisfies, in dimension three,
\[
 \partial_\tau l+\Delta l\le \frac{3/2-l}{\tau}.
\]
At an arbitrary endpoint of a minimizing $\mathcal L$-geodesic and for every
$\eta>0$, there is a smooth upper support $l_\eta$ which agrees with $l$ at
that endpoint and satisfies
\[
 \partial_\tau l_\eta+\Delta l_\eta
 \le \frac{3/2-l_\eta}{\tau}+\eta.
\]
\end{lemma}

\begin{proof}
\uses{def:reduced-length,def:ricci-flow-with-surgery}
For a variation field $Y$ along a minimizing curve $\gamma$, differentiate
the defining integral in \ref{def:reduced-length}.  Integration by parts,
using $\partial_tg=-2\operatorname{Ric}$, gives the endpoint term and the
$\mathcal L$-geodesic equation
\[
 \nabla_X X-\frac12\nabla R+2\operatorname{Ric}(X,\cdot)^*
 +\frac{1}{2\tau}X=0.
\]
Differentiate once more.  Trace the index form over fields with orthonormal
endpoint values and compare them with
$Y_i(s)=\sqrt{s/\tau}\,E_i(s)$, where $E_i$ is parallel after the Ricci
correction.  The boundary terms combine with the first-variation identities
for $l=L/(2\sqrt\tau)$, and the traced inequality is exactly
$\partial_\tau l+\Delta l\le(3/2-l)/\tau$.

At a cut endpoint, stop a fixed minimizing curve at a time
$0<\tau_0<\tau$.  In a neighborhood of its remaining segment, the reduced
exponential map is nonsingular.  Add the fixed $\mathcal L$-length of the
initial segment to the smooth terminal value function.  The resulting
function is at least $l$, equals it at the chosen endpoint, and obeys the
same second-variation calculation with an error tending to zero as
$\tau_0\downarrow0$.  Choosing $\tau_0$ so that this error is below $\eta$
gives the asserted support.
\end{proof}

\begin{lemma}[Minimum reduced-length comparison]
\label{lem:minimum-reduced-length-comparison}
\dcref{MT:ch7:7.10,MT:ch7:7.11,MT:ch7:7.12}
\source{morgan-tian}
\notready
Suppose that on every slice of a smooth backward-time interval the minimum
\[
 m(\tau)=\min_y l_x(y,\tau)
\]
is attained, and that the union of the minimizing sets over compact
subintervals is compact.  Then $m$ is continuous and
\[
 D^+m(\tau)\le\frac{3/2-m(\tau)}{\tau}.
\]
If $\limsup_{\tau\downarrow0}m(\tau)\le3/2$, then
$m(\tau)\le3/2$ throughout the interval.
\end{lemma}

\begin{proof}
\uses{lem:reduced-length-support-inequality,lem:forward-dini-mean-value}
Compactness of the minimizing set makes the upper supports in
\ref{lem:reduced-length-support-inequality} uniform for nearby slices.
Transport a minimizer by the time vector field for the upper estimate and a
nearby minimizer backward for the lower estimate.  The two estimates prove
continuity and, after division by the time increment and passage to the
limsup, the displayed Dini inequality.

Put $q=m-3/2$.  On every compact interval bounded away from zero,
$D^+q\le-q/\tau$.  If $q$ became positive, apply
\ref{lem:forward-dini-mean-value} to $e^{\int d\tau/\tau}q=\tau q$ between
the last zero and a later positive value.  The product has nonpositive upper
right Dini derivative but would have to increase, a contradiction.  Letting
the left endpoint decrease to zero proves $m\le3/2$.
\end{proof}

\begin{theorem}[Regularity of the reduced exponential map]
\label{thm:reduced-exponential-regularity}
\dcref{MT:ch6:6.18,MT:ch6:6.28,MT:ch6:6.56,MT:ch6:6.67}
\source{morgan-tian}
\notready
Assume that minimizing $\mathcal L$-geodesics to an open set remain in a
smooth region on which $|\operatorname{Ric}|$ and $|\nabla R|$ are bounded.
Then $\mathcal D_x^\tau$ is open,
\[
 \mathcal L\exp_x^\tau\colon\mathcal D_x^\tau
 \longrightarrow\Omega_x^\tau
\]
is a diffeomorphism, and $l_x(\cdot,\tau)$ is smooth on
$\Omega_x^\tau$.  The complement of $\Omega_x^\tau$ in the selected open
set has measure zero.
\end{theorem}

\begin{proof}
\uses{def:reduced-length,def:reduced-exponential-map,lem:reduced-length-support-inequality}
The $\mathcal L$-geodesic equation is a smooth ODE after the change of
variable $s=\sqrt\tau$.  Its derivative with respect to the initial vector is
the $\mathcal L$-Jacobi equation.  Absence of a Jacobi field vanishing at the
two endpoints makes this derivative invertible, so the inverse-function
theorem gives local smoothness of $\mathcal L\exp_x^\tau$.  Uniqueness of the
minimizer makes the local inverses agree, proving the diffeomorphism and the
smoothness of $l$.

The curvature bounds and the first-variation formula give a local Lipschitz
bound for $l$.  An endpoint with two distinct minimizing geodesics is a point
where $l$ has two different endpoint gradients, hence is not differentiable;
these points form a null set by Rademacher's theorem.  An endpoint with a
conjugate minimizer is a critical value of the smooth reduced exponential map
and therefore forms a null set by Sard's theorem.  Every point outside the
injectivity domain is of one of these two types, which proves the last claim.
\end{proof}

\begin{lemma}[Reduced-density monotonicity]
\label{lem:reduced-density-monotonicity}
\dcref{MT:ch6:6.50,MT:ch6:6.80,MT:ch6:6.81}
\source{morgan-tian}
\notready
Let $v\in\mathcal D_x^\tau$ and let $J(v,\tau)$ be the Jacobian of the
reduced exponential map.  Along the corresponding minimizing
$\mathcal L$-geodesic,
\[
 \frac d{d\tau}\left((4\pi\tau)^{-3/2}e^{-l}J\right)\le0.
\]
Consequently the reduced volume is nonincreasing in backward time as long as
all minimizing curves under consideration remain in one smooth domain.
\end{lemma}

\begin{proof}
\uses{def:reduced-volume,def:reduced-exponential-map,
  prop:reduced-injectivity-set-monotonicity,
  lem:reduced-length-support-inequality,thm:reduced-exponential-regularity}
Choose three $\mathcal L$-Jacobi fields whose endpoint values are an
orthonormal basis.  The logarithmic derivative of their volume is
$\partial_\tau\log J$.  The traced second-variation estimate used in
\ref{lem:reduced-length-support-inequality}, together with the endpoint
identities for $\partial_\tau l$, gives
\[
 \partial_\tau\log J-\partial_\tau l-\frac{3}{2\tau}\le0.
\]
This is the logarithmic derivative of the displayed density.  By
\ref{thm:reduced-exponential-regularity}, the injectivity domain has full
measure, so change of variables in \ref{def:reduced-volume} and monotone
convergence give the reduced-volume assertion without a cut-locus boundary
term.
\end{proof}

\begin{lemma}[Small reduced-exponential image estimate]
\label{lem:generalized-noncollapsing-small-image}
\dcref{MT:ch8:8.1}
\source{morgan-tian}
\uses{def:reduced-length,def:reduced-exponential-map,
  thm:local-ricci-flow-derivative-estimates,lem:reduced-length-support-inequality}
\notready
Let $(\mathcal M,G)$ be a generalized three-dimensional Ricci flow, let
$x\in M_T$, and let $0<\tau_0\le\bar\tau_0$, $0<r\le\sqrt{\tau_0}$.
Assume that $B_{g(T)}(x,r)$ has compact closure, admits a
time-vector-field-compatible embedding over $[T-r^2,T]$, and satisfies
$|\operatorname{Rm}|\le r^{-2}$ on that neighborhood.  Let
$\widetilde W\subset\widetilde{\mathcal U}_x(\tau_0)$ lie in the
$\mathcal L$-injectivity domain, with every associated geodesic having
reduced length at most $l_0$ on $[0,\tau_0]$, and put
\[
 \varepsilon=\frac{\operatorname{Vol}_{g(T)}B_{g(T)}(x,r)^{1/3}}{r},
 \qquad \tau_1=\varepsilon r^2,
 \qquad \widetilde W_{\rm sm}=
 \widetilde W\cap\{|v|\le(8\sqrt\varepsilon)^{-1}\}.
\]
There is $\varepsilon_0=\varepsilon_0(\bar\tau_0)>0$ such that, when
$\varepsilon\le\varepsilon_0$, the reduced volume of
$W_{\rm sm}(\tau_1)=\mathcal L\exp_x^{\tau_1}(\widetilde W_{\rm sm})$ satisfies
\[
 \widetilde V_x(W_{\rm sm}(\tau_1),\tau_1)
 \le 2\varepsilon^{3/2}.
\]
\end{lemma}

\begin{proof}
Rescale by $r^{-2}$.  The curvature bound and the local derivative estimate
give, on $B_{g(T)}(x,r/2)$ and $0\le\tau\le\tau_1$,
\[
 |\nabla R|\le C r^{-3},\qquad
 (1-C\varepsilon)g(T)\le g(T-\tau)\le(1+C\varepsilon)g(T).
\]
The reduced geodesic equation and the support inequality in
\ref{lem:reduced-length-support-inequality} imply, by a first-exit argument,
that every geodesic with $|v|\le(8\sqrt\varepsilon)^{-1}$ stays in
$B_{g(T)}(x,r/2)$ up to $\tau_1$.  Along it, $R\ge-6r^{-2}$, and hence
\[
 l_x(\gamma(\tau_1),\tau_1)\ge-2\varepsilon.
\]
The change-of-variables formula for the reduced exponential map and the
metric distortion estimate therefore give
\[
 \widetilde V_x(W_{\rm sm},\tau_1)
 \le e^{2\varepsilon}(1+C\varepsilon)\varepsilon^{-3/2}r^{-3}
       \operatorname{Vol}_{g(T)}B_{g(T)}(x,r).
\]
Since the last volume equals $\varepsilon^3r^3$, choosing
$\varepsilon_0$ small gives the stated factor $2$.  This is the small-vector
estimate in the proof of Morgan--Tian Theorem 8.1.
\end{proof}

\begin{lemma}[Large reduced-exponential tail estimate]
\label{lem:generalized-noncollapsing-large-tail}
\dcref{MT:ch8:8.1}
\source{morgan-tian}
\uses{def:reduced-exponential-map,prop:reduced-injectivity-set-monotonicity,
  def:reduced-volume,lem:reduced-density-monotonicity,
  thm:reduced-exponential-regularity}
\notready
Let $(\mathcal M,G)$ be a generalized three-dimensional Ricci flow with
the same compact-closure, compatible-embedding, and curvature hypotheses
as in the preceding lemma.  Let $x\in M_T$, let
$0<\tau_0\le\bar\tau_0$, $0<r\le\sqrt{\tau_0}$, and let
$\widetilde W\subset\widetilde{\mathcal U}_x(\tau_0)$ lie in the
$\mathcal L$-injectivity domain.  Assume moreover that there is an open
smooth spacetime set $\mathcal O$ containing $x$ and all
$\mathcal L\exp_x^\tau(\widetilde W)$ for $0<\tau\le\tau_0$, and that every
minimizing $\mathcal L$-geodesic from $x$ to those images stays in
$\mathcal O\cup\{x\}$.  Put
$\displaystyle \varepsilon=\operatorname{Vol}_{g(T)}B_{g(T)}(x,r)^{1/3}/r$,
and put
$\tau_1=\varepsilon r^2$ and
$\widetilde W_{\rm lg}=
\widetilde W\setminus\{|v|\le(8\sqrt\varepsilon)^{-1}\}$.  There is a
universal $\varepsilon_1>0$ such that, for
$0<\varepsilon\le\varepsilon_1$,
\[
 \widetilde V_x(\mathcal L\exp_x^{\tau_1}(\widetilde W_{\rm lg}),\tau_1)
 \le \varepsilon^{3/2}.
\]
\end{lemma}

\begin{proof}
The reduced-Jacobian monotonicity in
\ref{lem:reduced-density-monotonicity}, compared with the Euclidean
short-time asymptotics of the reduced exponential map, gives the pointwise
bound
\[
 \tau_1^{-3/2}e^{-\widetilde l(v,\tau_1)}J(v,\tau_1)
 \le 2^3 e^{-|v|^2}
\]
on the injectivity domain, after decreasing $\varepsilon_1$.  The regular
domain has full measure by the reduced-exponential regularity theorem, so no
cut-locus term is introduced.  Integrating over
$|v|\ge(8\sqrt\varepsilon)^{-1}$ and comparing a Euclidean ball with an
inscribed cube gives
\[
 \int_{|v|\ge(8\sqrt\varepsilon)^{-1}}2^3e^{-|v|^2}\,dv
 \le (4\pi)^{3/2}\frac32e^{-1/(192\varepsilon)}
 \le\varepsilon^{3/2}
\]
for sufficiently small $\varepsilon$.  This is the large-vector Gaussian
tail estimate in Morgan--Tian Theorem 8.1.
\end{proof}

\begin{theorem}[Imported interface: generalized-flow noncollapsing]
\label{thm:generalized-flow-noncollapsing-interface}
\dcref{MT:ch8:8.1}
\source{morgan-tian}
\notready
Fix $\bar\tau_0,l_0,V>0$.  There is
$\kappa=\kappa(\bar\tau_0,V,l_0)>0$ with the following property.  Let
$(\mathcal M,G)$ be a generalized three-dimensional Ricci flow, let
$x\in M_T$, and suppose that $0<\tau_0\le\bar\tau_0$ and
$0<r\le\sqrt{\tau_0}$.  Assume:
\begin{enumerate}
\item $B_{g(T)}(x,r)$ has compact closure in $M_T$;
\item the ball admits a time-vector-field-compatible embedding over
      $[T-r^2,T]$;
\item $|\operatorname{Rm}|\le r^{-2}$ on that parabolic neighborhood;
\item there is an open set
      $\widetilde W\subset\widetilde{\mathcal U}_x(\tau_0)\subset T_xM_T$
      in the $\mathcal L$-injectivity domain, together with an open smooth
      spacetime set $\mathcal O$ containing $x$ and every
      $W(\tau)=\mathcal L\exp_x^\tau(\widetilde W)$ for
      $0<\tau\le\tau_0$, such that every minimizing $\mathcal L$-geodesic
      from $x$ to $W(\tau)$ stays in $\mathcal O\cup\{x\}$ and every
      associated geodesic on $[0,\tau_0]$ has reduced length at most $l_0$;
\item for
      $W(\tau)=\mathcal L\exp_x^\tau(\widetilde W)$, one has
      $\operatorname{Vol}_{g(T-\tau_0)}W(\tau_0)\ge V$.
\end{enumerate}
Then
\[
 \operatorname{Vol}_{g(T)}B_{g(T)}(x,r)\ge\kappa r^3.
\]
\end{theorem}

\begin{proof}
\uses{def:reduced-volume,def:reduced-exponential-map,
  prop:reduced-injectivity-set-monotonicity,lem:reduced-density-monotonicity,
  thm:reduced-exponential-regularity,lem:generalized-noncollapsing-small-image,
  lem:generalized-noncollapsing-large-tail}
The reduced-volume contribution at the earlier time is bounded below directly
from hypothesis (5) and the reduced-length bound in (4):
\[
 \widetilde V_x(W(\tau_0),\tau_0)
 \ge (4\pi)^{-3/2}\bar\tau_0^{-3/2}Ve^{-l_0}=:c_0>0.
\]
By the star-shaped injectivity property in
\ref{prop:reduced-injectivity-set-monotonicity}, the same initial-vector set is injective
at every earlier time.  The common-domain hypothesis in
\ref{def:reduced-volume}, together with
\ref{lem:reduced-density-monotonicity}, therefore gives the same lower bound
at every $\tau\le\tau_0$, in particular at $\tau_1=\varepsilon r^2$ when
$\varepsilon\le1$.  Split the initial-vector domain into the small and large
sets in the two preceding lemmas.  Their estimates give
\[
 \widetilde V_x(W(\tau_1),\tau_1)
 \le3\varepsilon^{3/2}
\]
whenever $\varepsilon$ is below their common threshold.  Therefore either
$\varepsilon$ exceeds that threshold or
$\varepsilon\ge(c_0/3)^{2/3}$.  Since
$\operatorname{Vol}_{g(T)}B_{g(T)}(x,r)=\varepsilon^3r^3$, set
$\varepsilon_*:=\min\{\varepsilon_0,\varepsilon_1,1\}$ and take
\[
 \kappa=\min\{\varepsilon_*^3,(c_0/3)^2\}
\]
and this choice proves the claim.  The embedding and curvature hypotheses are
used in the small-image lemma.  The star-shaped injectivity and common-domain
clauses justify the comparison across times and the change of variables;
compact closure removes the cut-locus boundary term on the selected domain.
\end{proof}

\section{Local Curvature and Neck Estimates}

\begin{theorem}[Strong maximum principle for curvature tensors]
\label{thm:strong-curvature-tensor-maximum-principle}
\dcref{MT:ch4:4.17,MT:ch4:4.18}
\source{morgan-tian}
\uses{thm:time-slice-maximum-comparison}
\notready
Let $\overline U$ be a compact, connected, smooth codimension-zero submanifold
of a Riemannian manifold $M$, possibly with boundary, and let $\mathcal V$ be
a tensor bundle over $M$.  Let $\mathcal T(x,t)$ be a smooth family of
sections on $\overline U\times[0,T]$ satisfying
\[
 \partial_t\mathcal T=\Delta\mathcal T+\Psi(\mathcal T),
\]
where $\Psi$ is a smooth fiberwise vector field.  Suppose there is a
continuous, locally Lipschitz, fiberwise \emph{concave} function
$s\colon\mathcal V\to\mathbb R$ (equivalently, all of its superlevel sets are
convex) which is invariant under parallel translation and for which, whenever
$s_x(A)\ge0$, the vector $\Psi(A)$ lies in the tangent cone at $A$ to the
convex superlevel set $\{B\in V_x:s_x(B)\ge s_x(A)\}$.  If
$s(\mathcal T(x,0))\ge0$ on $\overline U$ and
$s(\mathcal T(x,t))\ge0$ on $\partial\overline U\times[0,T]$, and if
$s(\mathcal T(x_0,0))>0$ for some $x_0\in\operatorname{int}\overline U$,
then $s(\mathcal T(x,t))>0$ on
$\operatorname{int}\overline U\times(0,T]$.
\end{theorem}

\begin{proof}
Choose smooth $h(x,0)$ with $h=0$ on $\partial\overline U$,
$0\le h(x,0)\le s(\mathcal T(x,0))$, and $h(x_0,0)>0$, and let $h$ solve
the Dirichlet heat equation on $\overline U$.  The scalar strong maximum
principle gives $h>0$ on the later interior.  In the direct-sum bundle
$\mathcal V\oplus\mathbb R$, set
\[
 Z_x=\{(A,b):s_x(A)\ge b\ge0\}.
\]
Concavity and parallel invariance make $Z=\bigcup_xZ_x$ a closed
parallel-invariant convex subset.
The pair $(\mathcal T,h)$ evolves by a heat equation plus the fiberwise vector
field $(\Psi(A),0)$, which is tangent to $Z$ along its boundary by the
displayed hypothesis.  The tensor maximum principle therefore gives
$s(\mathcal T)\ge h$ throughout the cylinder.  At a nondifferentiable point
of $s$, use an affine supporting function of the concave superlevel set; the
same first-contact argument applies and is the standard viscosity form of
Morgan--Tian's tensor maximum principle.  The strict positivity of $h$ proves
the conclusion.
\end{proof}

\begin{lemma}[The three-dimensional two-smallest-curvature cone is reaction invariant]
\label{lem:two-smallest-curvature-cone-invariant}
\dcref{MT:ch4:4.17,MT:ch4:4.18}
\source{morgan-tian}
\uses{thm:strong-curvature-tensor-maximum-principle}
\notready
Let $\lambda\ge\mu\ge\nu$ be the eigenvalues of a three-dimensional
curvature operator, and put $s=\mu+\nu$.  For the Ricci-flow reaction ODE
\[
 \dot\lambda=\lambda^2+\mu\nu,\qquad
 \dot\mu=\mu^2+\lambda\nu,\qquad
 \dot\nu=\nu^2+\lambda\mu,
\]
the closed cone $\{s\ge0\}$ is invariant.  More precisely, on its boundary
 $s=0$ one has $\dot s=\mu^2+\nu^2\ge0$.  The function $s$ is a concave,
parallel-invariant function of the curvature operator, so its superlevel set
has the tangent-cone property required by
\ref{thm:strong-curvature-tensor-maximum-principle}.
\end{lemma}

\begin{proof}
On $s=0$, the ordering gives $\mu=-\nu\ge0$.  Adding the second and third
reaction equations gives
\[
 \dot s=\mu^2+\nu^2+\lambda(\mu+\nu)=2\mu^2\ge0.
\]
Thus the reaction vector points into the superlevel cone at every boundary
point.  The sum of the two smallest eigenvalues is the minimum of the traces
on two-planes complementary to a unit eigenline, hence is concave and is
unchanged by parallel transport.  The displayed calculation is exactly the
tangent condition used in Morgan--Tian's tensor maximum-principle proof.
\end{proof}

\begin{lemma}[Propagation of a curvature-operator nullity]
\label{lem:curvature-nullity-propagation}
\dcref{MT:ch4:4.19,MT:ch4:4.20}
\source{morgan-tian}
\uses{thm:strong-curvature-tensor-maximum-principle,
  lem:two-smallest-curvature-cone-invariant}
\notready
Let $T>0$ and let $(U,g(t))$, $0\le t\le T$, be a connected
three-dimensional Ricci flow whose slices have nonnegative sectional
curvature (completeness is not assumed).  If $g(0)$ is nonflat and the
Ricci tensor at some $p\in U$ at time $T$ has a zero eigenvalue, then
\[
 \mu+\nu=0
 \quad\text{and hence}\quad
 \operatorname{spec}(\operatorname{Rm})=\{\lambda,0,0\},\qquad\lambda>0,
\]
at every point of $U$ and every $0<t\le T$, where
$\lambda\ge\mu\ge\nu$ are the curvature-operator eigenvalues.  Thus the
Ricci tensor has a one-dimensional nullspace and its null line is defined
continuously (and smoothly after a local choice of sign).
\end{lemma}

\begin{proof}
For nonnegative sectional curvature, a zero Ricci eigenvalue at $(p,T)$
means that the two sectional curvatures containing the corresponding
direction vanish; equivalently $s(p,T):=\mu(p,T)+\nu(p,T)=0$.  We show
that $s$ vanishes everywhere.  If $s(q,t)>0$ at some earlier point, shift
time so that this point is at time zero and choose a compact connected
codimension-zero domain $V\Subset U$ whose interior contains both $p$ and
$q$.  The cone lemma gives $s\ge0$ on $V$ at the shifted initial time and
on $\partial V$ for the whole interval.  Applying the tensor strong
maximum principle
\ref{thm:strong-curvature-tensor-maximum-principle} with $q$ as the
strict interior point gives $s(p,T)>0$, a contradiction.  This is the
compact-domain time-shift argument in Morgan--Tian; no compact exhaustion
or terminal-time minimum is being assumed.

Since $g(0)$ is nonflat and the sectional curvatures are nonnegative, some
point has positive scalar curvature at time zero.  The scalar evolution and
the same Dirichlet maximum-principle argument on a compact domain containing
that point and any prescribed later point give $R>0$ for every $t>0$.
With $s=0$ and $\lambda,\mu,\nu\ge0$, we have $\mu=\nu=0$ and
$\lambda=R>0$, proving the spectrum and the null-line assertion.
\end{proof}

\begin{theorem}[Splitting at a curvature null direction]
\label{thm:curvature-null-splitting}
\dcref{MT:ch4:4.19,MT:ch4:4.20,MT:ch4:4.21}
\source{morgan-tian}
\uses{lem:curvature-nullity-propagation,
  lem:two-smallest-curvature-cone-invariant,
  thm:strong-curvature-tensor-maximum-principle}
\notready
Let $T>0$ and let $(U,g(t))$, $0\le t\le T$, be a connected smooth
three-dimensional Ricci flow with nonnegative sectional curvature; the
slices need not be complete for the local assertion.  Assume that $g(0)$ is
not flat and that the Ricci tensor at some $p\in U$ at time $T$ has a zero
eigenvalue (equivalently, the curvature operator has a two-dimensional
kernel).  Then, for every $0<t\le T$, the metric is locally a Riemannian
product of a surface of positive curvature and a line, and the flow is
locally the product of the surface Ricci flow and the trivial line flow.  If
every slice is complete, a covering
$\widetilde U\to U$ carries the product flow
$(N,h(t))\times(\mathbb R,ds^2)$ for $0\le t\le T$; deck
transformations preserve the parallel line distribution.  Without the
nonflatness assumption, the separate scalar clause is: if
$R(p,T)=0$, then the whole connected flow on $U\times[0,T]$ is flat.
\end{theorem}

\begin{proof}
\uses{lem:curvature-nullity-propagation,
  lem:two-smallest-curvature-cone-invariant,
  thm:strong-curvature-tensor-maximum-principle}
If $R(p,T)=0$ and the flow were nonflat, choose an earlier point with
positive scalar curvature, shift it to time zero, and choose a compact
connected codimension-zero domain whose interior contains that point and
$p$.  The scalar equation
$(\partial_t-\Delta)R=2|\operatorname{Ric}|^2\ge0$, with the boundary values
nonnegative, and the Dirichlet heat barrier give $R(p,T)>0$, a contradiction.
Hence $R\equiv0$; nonnegative sectional curvature then makes the whole
curvature tensor vanish.  This is the source's first strong-maximum
principle clause and uses no unsupported exhaustion limit.

For the nonflat case, apply the nullity-propagation lemma
\ref{lem:curvature-nullity-propagation}.  It gives the spectrum
$\{\lambda,0,0\}$ with $\lambda>0$ for $0<t\le T$ and supplies the
one-dimensional Ricci-null line distribution used below.

The two-dimensional curvature-operator kernel is dual to a tangent line
distribution $\mathcal D$.  Locally choose its unit generator $V$.  Along
any curve, parallel-extend $V$ and an arbitrary $W$.  Nonnegative sectional
curvature makes
$\operatorname{Rm}(\widetilde V,\widetilde W,\widetilde V,\widetilde W)$
nonnegative and zero at the base point, so its first variation vanishes.
Polarization gives $(\nabla\operatorname{Rm})(V,\cdot,\cdot,\cdot)=0$.
Differentiating
$\operatorname{Rm}(V,W_1,W_2,W_3)=0$ along three parallel fields shows that
$\nabla_XV$ lies in the line $\mathcal D$; unit length makes it orthogonal
to $V$, hence $\nabla_XV=0$.  Therefore the metric is locally a surface
of positive curvature times a line.  The curvature evolution preserves the
product tensor, and differentiating the kernel identity in the
time-dependent orthonormal gauge gives $\partial_tV=0$, so the splitting is
time-independent.

When the slices are complete, the complete-case globalization in the cited
Morgan--Tian corollary lifts this parallel line field to a product cover
$N\times\mathbb R$; deck transformations preserve the line field.  This is
the only global input beyond the local tensor argument and is included in
the source contract \dcref{MT:ch4:4.21}.
\end{proof}

\paragraph{Rank-one convention.}
The source's complete-case corollary is sometimes printed with the phrase
``$\operatorname{Rm}$ has a zero eigenvalue.''  The local proof uses the
stronger rank-one condition above, namely a two-dimensional kernel (or a
zero Ricci eigenvalue).  A single zero sectional eigenvalue does not imply
this condition, so no splitting conclusion is asserted for that weaker
literal reading.

\begin{theorem}[Ancient scalar-curvature monotonicity]
\label{thm:ancient-scalar-curvature-monotonicity}
\dcref{MT:ch4:4.37,MT:ch4:4.39}
\source{morgan-tian}
\notready
Let $(M,g(t))$, $-\infty<t\le0$, be a complete Ricci flow with bounded
nonnegative curvature operator on each time slice.  Then
\[
 \partial_tR(x,t)\ge0
\]
at every spacetime point.
\end{theorem}

\begin{proof}
Hamilton's differential Harnack inequality, applied on $(T_0,0]$ with zero
vector field, gives
\[
 \partial_tR+\frac{R}{t-T_0}\ge0.
\]
For fixed $(x,t)$ the curvature is finite and nonnegative.  Letting
$T_0\to-\infty$ makes the second term tend to zero and proves the assertion.
\end{proof}

\begin{lemma}[Finite-interval Harnack and distance control]
\label{lem:finite-interval-harnack-distance-control}
\dcref{MT:ch11:11.12}
\source{morgan-tian}
\notready
Let $0<T<\infty$ and let $(M,g(t))$, $-T<t\le0$, be a complete
three-dimensional Ricci flow with nonnegative curvature and curvature locally
bounded in time.  If
\[
 Q=\sup_{x\in M}R_{g(0)}(x)<\infty,
\]
then for every $x\in M$ and $-T<t\le0$,
\[
 R_{g(t)}(x)\le Q\frac{T}{T+t},
\]
and for all $x,y\in M$,
\[
 d_{g(t)}(x,y)\le d_{g(0)}(x,y)+16\sqrt{\frac Q3}\,T.
\]
The estimates also hold on compact exhaustions when the flow is initially
given only as a pointed geometric limit.
\end{lemma}

\begin{proof}
On a compact exhaustion, Hamilton's differential Harnack inequality gives
\(\partial_tR\ge-R/(t+T)\).  Integrating from $t$ to $0$ yields the
displayed scalar bound.  Nonnegative curvature in dimension three gives
\(\operatorname{Ric}\le (n-1)R/2\); the metric-distortion inequality and
integration of the resulting square-root bound give the constant
\(16\sqrt{Q/3}\,T\).  The estimates are independent of the exhaustion, so
local smooth convergence and passage to the limit give the final assertion.
This is Morgan--Tian Claim 11.12, with the finite-time Harnack and distance
steps stated explicitly.
\end{proof}

\begin{theorem}[Weak superharmonicity of Busemann functions]
\label{thm:busemann-function-superharmonicity}
\dcref{MT:ch2:2.3}
\source{morgan-tian}
\notready
Let $(M^n,g)$ be complete with nonnegative Ricci curvature, and let $B$ be
the Busemann function of a ray.  Then $\Delta B\le0$ weakly: for every
nonnegative compactly supported smooth function $\varphi$,
\[
 \int_M B\,\Delta\varphi\,d\mu\le0.
\]
\end{theorem}

\begin{proof}
For the ray $\gamma$, put
$B_T(x)=d(x,\gamma(T))-T$.  Laplacian comparison away from the cut locus,
interpreted by upper barriers, gives
\[
 \int_M B_T\Delta\varphi\,d\mu
 \le \int_M\frac{n-1}{d(x,\gamma(T))}\varphi\,d\mu.
\]
If the support of $\varphi$ lies in the $D$-ball about $\gamma(0)$, the
denominator is at least $T-D$.  The right side therefore tends to zero.
The functions $B_T$ converge locally uniformly to $B$ by monotonicity and
the triangle inequality, so passage to the limit proves the weak inequality.
\end{proof}

\begin{lemma}[Long minimizing geodesics align with a neck axis]
\label{lem:epsilon-neck-geodesic-axis-alignment}
\dcref{MT:ch19:19.4}
\source{morgan-tian}
\uses{def:epsilon-neck-and-scale}
\notready
For every $\alpha>0$ there is $\epsilon_0(\alpha)>0$ such that, in a
normalized $\epsilon$-neck with $\epsilon\le\epsilon_0$, every minimizing
geodesic of length greater than $\epsilon^{-1}/100$ which crosses in the
positive axial direction has tangent within $\alpha$ of the axial unit
vector at every point of the geodesic contained in the neck.
\end{lemma}

\begin{proof}
Suppose not and take $\epsilon_j\downarrow0$.  Pull the metrics and geodesics
back to the standard cylinders.  On every fixed compact subcylinder the
metrics converge in $C^2$ to the product metric and the geodesics converge
to a minimizing product geodesic.  Its axial displacement tends to infinity,
whereas its spherical displacement is at most the diameter of the round
two-sphere.  Recenter at any chosen point of misalignment.  On the resulting
full- or half-cylinder limit, the conserved product momenta force its
spherical speed to zero and its positive axial speed to one.  Smooth
convergence of the geodesic equation gives the same estimate for the
approximating tangent at the recentered point, a contradiction.
\end{proof}

\section{Compactness and Noncollapsing}

\begin{theorem}[Pointed Riemannian compactness]
\label{thm:pointed-riemannian-compactness}
\dcref{MT:ch5:5.6,MT:ch5:5.9,MT:ch5:5.10}
\source{morgan-tian}
\notready
Let $(M_j,g_j,x_j)$ be complete pointed Riemannian manifolds.  Suppose that
on every fixed based ball all covariant derivatives of curvature are
uniformly bounded and that for some $r_0,v_0>0$,
\[
 \operatorname{Vol}_{g_j}B(x_j,r_0)\ge v_0.
\]
Then a subsequence converges smoothly on compact subsets to a complete
pointed Riemannian manifold.
\end{theorem}

\begin{proof}
\uses{thm:local-volume-injectivity-radius-control}
The base-ball curvature and volume bounds give a positive lower bound for
the injectivity radius at $x_j$ by
\ref{thm:local-volume-injectivity-radius-control}.  Bishop--Gromov comparison
and the curvature bounds propagate the same type of volume lower bound to
centers in each fixed based ball, so the same theorem gives a uniform
injectivity-radius bound there.

Exponential charts of a fixed radius therefore cover every such ball with a
uniform number of charts.  In harmonic refinements of these charts, the
curvature-derivative hypotheses give uniform $C^k$ bounds for the metric
coefficients.  Arzela--Ascoli and a diagonal choice in the radius and $k$
give compatible limiting charts and smooth transition maps.  They assemble
to the pointed smooth limit.  Every finite-length limiting geodesic remains
inside one ball of the exhaustion and is the limit of extendible geodesics
in $M_j$; hence it extends.  Hopf--Rinow gives completeness.
\end{proof}

\begin{theorem}[Pointed Ricci-flow compactness]
\label{thm:hamilton-compactness-generalized-flows}
\dcref{MT:ch5:5.14,MT:ch5:5.15}
\source{morgan-tian}
\notready
Let $(\mathcal M_j,G_j,x_j)$ be based generalized Ricci flows with base time
zero, defined on a common open time interval $I$.  Suppose that for every
$A<\infty$ and compact $J\Subset I$ containing zero, all sufficiently large
$j$ admit compatible embeddings
\[
 B_{g_j(0)}(x_j,A)\times J\longrightarrow\mathcal M_j
\]
with compact zero-time balls and a uniform curvature bound on the images.
Assume that for some $r_0,\kappa>0$,
\[
 \operatorname{Vol}_{g_j(0)}B(x_j,r_0)\ge\kappa r_0^3.
\]
After passing to a subsequence, the flows converge smoothly on compact
spacetime subsets.  The limit is a pointed Ricci flow
\[
 (M_\infty,g_\infty(t),x_\infty)
\]
on $I$, and its time-zero slice is complete.  If
$\operatorname{Ric}_{g_\infty(t)}\ge0$ for $t\le0$, then every nonpositive
time slice of the limit is complete.
\end{theorem}

\begin{proof}
\uses{def:ricci-flow-with-surgery,def:surgery-flow-noncollapsing,thm:local-ricci-flow-derivative-estimates,thm:pointed-riemannian-compactness}
The curvature bound and the Ricci-flow equation give uniform bounds for all
covariant derivatives of curvature on smaller compact cylinders by
\ref{thm:local-ricci-flow-derivative-estimates}.  The base-ball volume lower
bound and curvature bounds put the zero-time slices under
\ref{thm:pointed-riemannian-compactness}, producing a complete smooth
zero-time limit and compatible exhaustion charts.

Exhaust each zero-time slice by the compact balls and exhaust $I$ by compact
intervals.  Repeating the construction and diagonalizing produces compatible
limit charts; their transition maps agree because they are limits of the
original embeddings.  The limiting metrics satisfy the Ricci-flow equation
by smooth convergence, and pointed Riemannian compactness already gives
completeness at time zero.

If the limiting Ricci curvature is nonnegative and $t\le0$, then
$g_\infty(t)\ge g_\infty(0)$ as quadratic forms.  A finite-length
$g_\infty(t)$-curve escaping every compact set would therefore have finite
$g_\infty(0)$-length, contradicting completeness of the time-zero slice.
Thus every nonpositive slice is complete under the additional curvature
hypothesis.  Local based-ball bounds alone make no completeness assertion at
other times.
\end{proof}

\begin{lemma}[Nonflat terminal cones are excluded]
\label{lem:nonflat-terminal-cone-exclusion}
\dcref{MT:ch4:4.22}
\source{morgan-tian}
\notready
Let $T>0$ and let $(U,g(t))$, $0\le t\le T$, be a connected
three-dimensional Ricci flow (the slices need not be complete) with
nonnegative sectional curvature.  Suppose that, for a Riemannian
two-manifold $(N,h)$, the terminal slice $(U,g(T))$ is isometric to a
nonempty open subset of the open cone
\[
 (N\times(0,\infty),\;dr^2+r^2h).
\]
Then $(U,g(t))$ is flat for every $t\in[0,T]$.
\end{lemma}

\begin{proof}
\uses{thm:curvature-null-splitting}
This is the no-cone proposition of Morgan--Tian, with its hypotheses and
time interval stated explicitly.  For completeness, the source argument is
as follows.  If the terminal cone is flat, the vanishing-scalar-curvature
strong maximum principle gives flatness for every earlier time.  If it is
nonflat, the cone curvature calculation gives a two-dimensional curvature
kernel and a nonzero eigenvalue that scales as $r^{-2}$ along cone rays.
The local product splitting in
\ref{thm:curvature-null-splitting} makes the null line parallel and hence
forces that eigenvalue to be constant along its line, a contradiction.
Thus the terminal slice is flat, and the scalar clause in
\ref{thm:curvature-null-splitting} propagates flatness to all
$t\in[0,T]$.  No completion of the cone vertex or removable-singularity
claim is used.
\end{proof}

\begin{lemma}[Normalized standard-cap reduced-length barrier]
\label{lem:standard-cap-reduced-length-barrier}
\dcref{MT:ch16:16.5,MT:ch16:16.8,MT:ch16:16.13}
\source{morgan-tian}
\uses{def:standard-surgery-cap-model,lem:local-surgery-standard-cap-convergence,
  thm:hamilton-compactness-generalized-flows}
\notready
For every $L_0<\infty$ and every fixed surgery-cap model, there are a
normalized inner cap neighborhood, an outer annulus, and a time
$0<\theta<1$ such that every backward reduced path which enters the inner
neighborhood from the outer side before normalized time $\theta$ has
$\mathcal L_+\ge L_0$.  The same lower bound holds for a path ending at the
normalized exposed frontier.  Standard-cap convergence transfers the barrier
uniformly to all sufficiently small surgery controls, with the terminal
controlled ball outside the outer annulus.
\end{lemma}

\begin{proof}
Rescale by the cap scale.  On each fixed normalized ball, the inserted metrics
and flows converge smoothly to the standard cap by
\ref{lem:local-surgery-standard-cap-convergence}; Hamilton compactness gives
the same convergence on compact subintervals $[0,\theta]$.  Choose the outer
annulus wide enough that the weighted Cauchy--Schwarz lower bound for crossing
it exceeds $L_0$.  A path remaining in the inner cap until time $\theta$ has
the complementary scalar-curvature lower bound from the standard solution;
choosing $\theta$ close enough to one makes this exceed $L_0$ as well.  Smooth
convergence transfers both alternatives, including the exposed-frontier
limit, uniformly once the surgery control is small.  This is the normalized
barrier estimate used in Morgan--Tian's Chapter~16 argument.
\end{proof}

\begin{lemma}[Short reduced paths avoid surgery caps]
\label{lem:short-reduced-paths-avoid-surgery-caps}
\dcref{MT:ch16:16.5,MT:ch16:16.8,MT:ch16:16.13,MT:ch16:16.15}
\source{morgan-tian}
\uses{def:controlled-ricci-flow-with-surgery,def:reduced-length,
  lem:standard-cap-reduced-length-barrier}
\notready
Fix a controlled stage with a finite list of surgery times in
$[T_{i-1},T]$, a terminal point $(x,T)$ with a controlled parabolic ball of
radius at least $r_{i+1}$, and $L_0<\infty$.  There is a
positive surgery-control bound, depending only on the stage data and $L_0$,
such that every backward path from $x$ with
\[
 \mathcal L_+(\gamma)=\int\sqrt\tau
   \bigl(\max(R,0)+|\dot\gamma|^2\bigr)\,d\tau<L_0
\]
misses a fixed open neighborhood of every surgery cap and every exposed
point in $[T_{i-1},T]$.  The finite list is part of the stage hypothesis; the
neighborhood and control bound are uniform over that list.
\end{lemma}

\begin{proof}
For $0<a<b$, Cauchy--Schwarz gives
\[
 \int_a^b\sqrt\tau|\dot\gamma|^2d\tau
 \ge\frac{\sqrt a\,d(\gamma(a),\gamma(b))^2}{b-a}.
\]
The normalized barrier in
\ref{lem:standard-cap-reduced-length-barrier} applies to every cap and exposed
frontier in the finite surgery list.  The terminal controlled ball is outside
the corresponding outer annuli, so a path from $(x,T)$ which entered one of
them would have reduced length at least $L_0$, contradicting the hypothesis.
The finitely many barriers can be intersected by taking the minimum of their
control bounds.  Thus the same fixed neighborhood size works for every cap and
exposed point in the stage.
\end{proof}

\begin{proposition}[Compact smooth domain for reduced minimizers]
\label{prop:compact-reduced-minimizing-domain}
\dcref{MT:ch16:16.4,MT:ch16:16.21,MT:ch16:16.24,MT:ch16:16.25}
\source{morgan-tian}
\notready
Under the hypotheses of \ref{lem:short-reduced-paths-avoid-surgery-caps},
there is an open subset $U$ of the smooth spacetime over
$[T_{i-1},T)$ such that:
\begin{enumerate}
\item every point of $U$ is joined to $x$ by a minimizing smooth
      $\mathcal L$-geodesic;
\item every slice $U_t$ is nonempty and
      $\min_{U_t}\mathcal L\le3\sqrt{T-t}$;
\item over every compact time interval, the slice minimizers form a compact
      subset of the smooth spacetime.
\end{enumerate}
\end{proposition}

\begin{proof}
\uses{lem:short-reduced-paths-avoid-surgery-caps,lem:minimum-reduced-length-comparison}
The scalar lower bound $R\ge-6$ gives
\[
 \mathcal L_+(\gamma)
 \le\mathcal L(\gamma)+4T_{i+1}^{3/2}.
\]
Remove the cap neighborhoods supplied by
\ref{lem:short-reduced-paths-avoid-surgery-caps}.  There are only finitely
many singular times on the compact interval already constructed, and every
slice is compact, so the complement $Y(L_0)$ is compact and smooth.  It
contains every path whose $\mathcal L$-length is below the chosen bound.

A minimizing sequence in $Y(L_0)$ has a uniformly bounded weighted kinetic
energy.  Away from backward time zero this gives equicontinuity by
Cauchy--Schwarz; inside the terminal controlled ball the curvature and metric
distortion estimates give the same conclusion up to zero.  Arzela--Ascoli
gives uniform convergence, weak $W^{1,2}$ compactness gives lower
semicontinuity of the kinetic term, and smooth convergence in $Y(L_0)$ gives
convergence of the scalar term.  The limit is therefore a smooth minimizing
$\mathcal L$-geodesic.

Let $U$ consist of endpoints joined to $x$ by paths below one fixed bound
larger than $8\sqrt{T_{i+1}}(1+T_{i+1})$.  The direct method gives the first
claim, and the same compact set with varying endpoint time gives the third.
Near backward time zero, the controlled terminal ball has minimum reduced
length at most $3/2$.  Apply
\ref{lem:minimum-reduced-length-comparison}; if the interval on which this
bound holds ended early, compactness of its minimizers and a short vertical
extension would continue it.  Thus the bound holds on every slice and gives
the second claim.
\end{proof}

\begin{lemma}[Positive-curvature components persist along smooth flow lines]
\label{lem:positive-component-persistence}
\dcref{MT:ch4:4.23,MT:ch16:16.1}
\source{morgan-tian}
\uses{def:ricci-flow-with-surgery,def:curvature-pinched-toward-positive}
\notready
Let $X_s$ be a compact connected component of a smooth three-dimensional
Ricci flow with strictly positive sectional curvature at a time $s$.  Until
the next surgery or extinction, its forward images remain compact and have
strictly positive sectional curvature.  Consequently, if a smooth backward
flow line joins a point $y\in X_s$ to a later point $x$, then the component
containing $x$ is a positive-sectional-curvature component.
\end{lemma}

\begin{proof}
The curvature-operator maximum principle in Morgan--Tian 4.23 preserves the
open positive-curvature cone on the smooth interval.  The time-vector-field
identification carries a connected component to its forward image until a
surgery boundary is met; a smooth flow line cannot cross that boundary.  The
last assertion follows.
\end{proof}

\begin{proposition}[Morgan--Tian stagewise reduced-exponential family]
\label{prop:stagewise-reduced-exponential-family}
\dcref{MT:ch16:16.1,MT:ch16:16.4,MT:ch16:16.15,MT:ch16:16.21,MT:ch16:16.26}
\source{morgan-tian}
\uses{def:controlled-ricci-flow-with-surgery,def:reduced-exponential-map,
  def:controlled-surgery-parameters,
  lem:short-reduced-paths-avoid-surgery-caps,prop:compact-reduced-minimizing-domain,
  thm:reduced-exponential-regularity,def:surgery-flow-noncollapsing,
  lem:positive-component-persistence,thm:local-ricci-flow-derivative-estimates}
\notready
Fix $i\ge1$ and the finite control data of
\ref{def:controlled-surgery-parameters}.  Let $T\in[T_i,T_{i+1})$ and let
$(x,T)\in\mathcal M_T^-$ be a smooth presurgery point in a component which is
not compact and positive sectional curvature.  Assume the restriction before
$T_i$ has the stage
$\kappa_i$-noncollapsing bound, the whole flow satisfies Assumptions (1)--(7),
curvature pinching, and the strong $(C,\epsilon)$-canonical-neighborhood
assumption with parameter $r_{i+1}$, and
\[
 \bar\delta(t)\le\mathfrak d_i(r_{i+1})
 \qquad(T_{i-1}\le t<T_{i+1}).
\]
Suppose that the fixed stagewise selector has supplied a value
$0<\mathfrak d_i(r_{i+1})<1/2$ and that, for some
$\rho\in[r_{i+1},\epsilon]$ with $T\ge\rho^2$, the compatible backward
parabolic neighborhood
$P(x,\rho,T,-\rho^2)$ exists and
$|\operatorname{Rm}|\le\rho^{-2}$ there.  Then there are a time
$\tau_*>0$, an open smooth domain $U\subset\mathbf t^{-1}[T_{i-1},T)$,
constants $l_0<\infty$, $r'_i>0$, $V_i=\kappa_i(r'_i)^3>0$, and an open set
$\widetilde W\subset\widetilde{\mathcal U}_x(\tau_*)$ in the reduced-
exponential injectivity domain such that:
\begin{enumerate}
\item $U_t$ is nonempty for every $t\in[T_{i-1},T)$, every point of $U$
      is joined to $(x,T)$ by a minimizing smooth $\mathcal L$-geodesic,
      and the slice minimum is at most $3\sqrt{T-t}$;
\item the set of slice minimizers is compact over every compact subinterval
      of $[T_{i-1},T)$;
\item there are $t''=T-\tau_*$ and $y\in U_{t''}$ with
      $R(y)<r_i^{-2}$ and a compatible parabolic ball
      $P(y,t'',r'_i,-(r'_i)^2)$ on which
      $|\operatorname{Rm}|\le(r'_i)^{-2}$ and
      $\operatorname{Vol}_{g(t'')}B(y,t'',r'_i)\ge V_i$;
\item the terminal ball $B_{g(T)}(x,\rho)$ has compact closure, embeds
      compatibly and unscathed over $[T-\rho^2,T]$, and has
      $|\operatorname{Rm}|\le\rho^{-2}$ there;
\item every reduced geodesic issued from $\widetilde W$ has reduced length at
      most $l_0$ on $[0,\tau_*]$, and
      $\operatorname{Vol}_{g(T-\tau_*)}
      \mathcal L\exp_x^{\tau_*}(\widetilde W)\ge V_i$.
\end{enumerate}
Here $r'_i$ and $l_0$ depend only on the fixed stage data (and not on $x$ or
$T$).
\end{proposition}

\begin{proof}
\uses{thm:generalized-flow-noncollapsing-interface}
Set $L=8\sqrt{T_{i+1}}(1+T_{i+1})$.  The cap-evolution estimates in
\ref{lem:short-reduced-paths-avoid-surgery-caps} and the compact minimizing
domain proposition (Morgan--Tian Proposition 16.4) give $U$, the slice
minimizers, and the $3\sqrt{T-t}$ bound.  Choose one minimizing curve to the
slice $T_{i-1}$.  If $R\ge r_i^{-2}$ on the interval between
$\max\{\epsilon^2,T-T_i\}$ and $T-T_{i-1}-\epsilon^2$, then the scalar part of
its reduced length, while $R\ge-6$ on the two end intervals, gives
\[
 \mathcal L\ge4\sqrt{T-T_{i-1}},
\]
contradicting the $3\sqrt{T-T_{i-1}}$ bound.  Thus it reaches a point
$y$ with $R(y)<r_i^{-2}$.

Put $\Delta=\min\{r_i^2/(16C),\epsilon^2\}$.  The local derivative and
curvature estimates give $C'<\infty$ on the backward ball about $y$ for time
$\Delta$; the cap-curvature lower bound, together with the displayed smallness
of $\bar\delta$, prevents that ball from meeting a surgery cap.  Metric
distortion supplies $\Delta_1\le\Delta/2$.  Taking
\[
 r'_i\le\min\{r_i/(32C),\Delta_1/2,(C')^{-1/2},\sqrt{\Delta_1}/2\}
\]
produces the parabolic ball in item (3).  The inductive noncollapsing estimate
gives its volume lower bound unless its component is positively curved; in
that exceptional case \ref{lem:positive-component-persistence} would make the
terminal component positive, contrary to the hypothesis.

For each $z$ in the earlier ball, join $y$ to $z$ by the short spatial geodesic
inside the controlled ball and concatenate it with the initial segment from
$x$.  The metric comparison gives $\mathcal L(\gamma_z)\le L/2$.
Proposition 16.4 then puts the minimizing replacement in the smooth domain.
The reduced exponential regularity theorem and the full-measure cut-locus
statement select an open full-measure subset
$B'\subset B(y,t'',r'_i)$ on which the reduced exponential is injective.  Set
$\widetilde W$ to be its inverse image.  The volume of its image is at least
$\kappa_i(r'_i)^3=V_i$, and the concatenation bound gives $l_0=L/2$.
This proves all five hypotheses of
\ref{thm:generalized-flow-noncollapsing-interface} explicitly.
\end{proof}

\begin{lemma}[Reduced-volume noncollapsing transfer]
\label{lem:reduced-volume-noncollapsing-transfer}
\dcref{MT:ch16:16.1,MT:ch16:16.4,MT:ch16:16.26,MT:ch6:6.80}
\source{morgan-tian}
\notready
Assume the stage-$i$ noncollapsing bound on earlier controlled slices and all
stage hypotheses of
\ref{prop:stagewise-reduced-exponential-family}.  For each selected terminal
centre choose $\tau_*$, $\widetilde W$, $l_0$, and $V_i$ from that proposition;
these data explicitly provide all five hypotheses of
\ref{thm:generalized-flow-noncollapsing-interface}.  There is then
$\kappa_{\mathcal L}>0$, depending only on the stage data, such that every
terminal parabolic ball of radius
$\rho\in[r_{i+1},\epsilon]$ and curvature at most $\rho^{-2}$ satisfies
\[
 \operatorname{Vol}B(x,T,\rho)\ge\kappa_{\mathcal L}\rho^3,
\]
for every selected terminal centre whose component does not have positive
sectional curvature.  This is deliberately the nonpositive branch; the
positive-component branch is supplied by the Kleiner--Lott estimate below.
\end{lemma}

\begin{proof}
\uses{def:controlled-ricci-flow-with-surgery,def:curvature-pinched-toward-positive,
  def:strong-canonical-neighborhood-control,prop:stagewise-reduced-exponential-family,
  thm:generalized-flow-noncollapsing-interface}
For each terminal centre, apply
\ref{prop:stagewise-reduced-exponential-family}.  Items (3)--(5) of that
proposition are respectively compact closure and compatible curvature control,
the uniform reduced-length bound, and the positive earlier image volume.
Thus all five hypotheses of
\ref{thm:generalized-flow-noncollapsing-interface} are present before it is
invoked, and its constant depends only on the fixed stage data.  The positive
sectional-curvature branch is intentionally excluded here and is handled by
the all-cases Kleiner--Lott estimate.
\end{proof}

\begin{lemma}[Kleiner--Lott all-cases noncollapsing estimate]
\label{lem:kleiner-lott-noncollapsing-estimate}
\dcref{KL:II.5:II.5.2}
\source{kleiner-lott}
\uses{def:normalized-riemannian-three-manifold,def:ricci-flow-with-cutoff,
  def:surgery-flow-noncollapsing,lem:reduced-volume-noncollapsing-transfer,
  thm:generalized-flow-noncollapsing-interface,lem:reduced-length-support-inequality,
  thm:hamilton-compactness-generalized-flows}
\notready
Fix $0<r_+\le r_-\le\epsilon$, $\kappa_->0$, and finite numbers
$0<E_-\le E$.  There are constants
\[
 \bar\delta=\bar\delta(r_-,r_+,\kappa_-,E_-,E)>0,
 \qquad
 \kappa_+=\kappa_+(r_-,\kappa_-,E_-,E)>0
\]
with the following property.  Let $a<b<c$, $b-a\ge E_-$, and
$c-a\le E$.  Let $(\mathcal M,G)$ be a three-dimensional Ricci flow with
Kleiner--Lott $(r,\delta)$-cutoff on an interval containing $[a,c)$, with
normalized compact initial data, $r\le\epsilon$, $r\ge r_-$ on $[a,b)$,
and $r\ge r_+$ on $[b,c)$.  If the flow is ordinarily
$\kappa_-$-noncollapsed below $\epsilon$ on $[a,b)$ and
$\delta(t)\le\bar\delta$ on $[a,c)$, and it satisfies the strong
canonical-neighborhood assumption with parameter $r_+$ on $[b,c)$, then it is ordinarily
$\kappa_+$-noncollapsed below $\epsilon$ on $[b,c)$ in the sense of
\ref{def:surgery-flow-noncollapsing}; at each test centre it gives either the
volume lower bound or the positive-sectional-curvature alternative in that
definition.
\end{lemma}

\begin{proof}
The canonical-neighborhood assumption on $[b,c)$ is used exactly at the
first-failure barrier: it keeps every competitive reduced-length curve away
from the surgery caps and supplies the neck/cap curvature and volume models.
The cited proof first treats a test ball of radius below
$\sqrt{E_-/3}$ by choosing a reduced-length minimizer back to time $a$.
The lower scalar-curvature bound and the $\mathcal L$ support inequality
force a point on the minimizer into the earlier interval with curvature below
$r_-^{-2}$.  The induction hypothesis gives a definite volume ball there;
the compact reduced-minimizing-domain and generalized-flow noncollapsing
calculation transfer that volume to the terminal ball.  The surgery-cap
barrier excludes a cap from the transferred parabolic region.

For the remaining radii, the canonical-neck and standard-cap estimates give a
uniform volume ratio directly whenever the centre is not in a positive
sectional-curvature component.  The source's Proposition II.5.2 explicitly
leaves the latter case as the positive-curvature alternative; no compactness
argument in this Blueprint upgrades it to a volume bound.  Taking the minimum
of the constants in the nonpositive branch gives $\kappa_+$, which depends
only on $(r_-,\kappa_-,E_-,E)$ and not on $r_+$.
\end{proof}

\begin{definition}[Stagewise surgery-control selector]
\label{def:stagewise-delta-selector}
\dcref{KL:II.5:II.5.2}
\source{kleiner-lott}
\notready
For fixed stage data $(i,\delta_i,r_i,\kappa_i,T_i)$, a
\emph{stagewise surgery-control selector} is a function
\[
 \mathfrak d_i:(0,r_i]\longrightarrow(0,\delta_i]
\]
such that the conditional extension conclusion below holds with the cutoff
bound $\bar\delta\le\mathfrak d_i(r)$ at radius $r$.  Whenever the
pointwise existence statement is used for a family of radii, one selector is
fixed once and for all by the axiom of choice; the notation
$\delta(r)$ in later first-failure arguments means this fixed
$\mathfrak d_i(r)$ for the relevant stage, not an unquantified new witness.
\end{definition}

\begin{theorem}[Conditional noncollapsing extension]
\label{thm:conditional-noncollapsing-extension}
\dcref{MT:ch16:16.1,KL:II.5:II.5.2}
\source{kleiner-lott}
\uses{def:normalized-riemannian-three-manifold,def:ricci-flow-with-cutoff,
  def:surgery-flow-noncollapsing,lem:kleiner-lott-noncollapsing-estimate,
  def:stagewise-delta-selector}
\notready
Fix $i\ge0$ and positive nonincreasing finite sequences
\[
\delta_0\ge\cdots\ge\delta_i,
\qquad \epsilon=r_0\ge\cdots\ge r_i,
\qquad \kappa_0\ge\cdots\ge\kappa_i.
\]
Fix stage times with $T_{-1}=0$ and explicit bounds
\[
 0<E_-\le T_i-T_{i-1},\qquad T_{i+1}-T_{i-1}\le E<\infty.
\]
There are $0<\kappa_{i+1}\le\kappa_i$ and a stagewise selector
$\mathfrak d_i:(0,r_i]\to(0,\delta_i]$ in the sense of
\ref{def:stagewise-delta-selector} such that, for every
$0<r_{i+1}\le r_i$, the following property holds with
$\delta(r_{i+1})=\mathfrak d_i(r_{i+1})$.  If a
three-dimensional Ricci flow with Kleiner--Lott $(r,\bar\delta)$-cutoff on
$[0,T)$, $T\in(T_i,T_{i+1}]$, has normalized compact initial data in the
sense of \ref{def:normalized-riemannian-three-manifold}, satisfies the
controlled-surgery topological assumptions, and is ordinarily
  $\kappa_i$-noncollapsed below $\epsilon$ on
  $[T_{i-1},T_i)$ (with $T_{-1}=0$), satisfies

  \[
   r(t)\ge r_i\quad(t\in[T_{i-1},T_i)),
   \qquad
   r(t)\ge r_{i+1}\quad(t\in[T_i,T)),
  \]
  and $r(t)\le\epsilon$ throughout $[T_{i-1},T)$,
  and has a positive nonincreasing
  surgery-control function satisfying
\[
 \bar\delta(t)\le\delta_j
 \quad\bigl(t\in[T_j,T_{j+1}),\ 0\le j\le i\bigr),
 \qquad
 \bar\delta(t)\le\delta(r_{i+1})
 \quad\bigl(t\in[T_{i-1},T_{i+1})\bigr),
\]
  where $T_{-1}=0$, and satisfies the strong canonical-neighborhood
  assumption with parameter $r_{i+1}$ on $[T_i,T)$, then every parabolic ball
  of radius at most
  $\epsilon$ satisfying the curvature hypothesis in
\ref{def:surgery-flow-noncollapsing} satisfies its source disjunction with
constant $\kappa_{i+1}$ for every centre and every component of the time slice:
either the ball has volume at least $\kappa_{i+1}\rho^3$, or its centre lies in
a positive-sectional-curvature component.  The constant is independent of the
particular flow.
\end{theorem}

\begin{proof}
\uses{def:surgery-flow-noncollapsing,def:controlled-ricci-flow-with-surgery,
  def:ricci-flow-with-cutoff,lem:kleiner-lott-noncollapsing-estimate}
Put $a=T_{i-1}$ (and $a=0$ when $i=0$), $b=T_i$, and $c=T$.  The
  hypotheses supply normalized initial data and the $(r,\delta)$-cutoff.  The
  finite stage interval gives fixed numbers $E_-$ and
$E$ with $b-a\ge E_-$ and $c-a\le E$; the induction hypothesis is ordinary
$\kappa_i$-noncollapsing on $[a,b)$, and the imposed bound on $\bar\delta$ is
exactly the cutoff bound required by
\ref{lem:kleiner-lott-noncollapsing-estimate}.  Apply that lemma with
$r_-=r_i$ and $r_+=r_{i+1}$.  It returns a positive constant independent of
the flow and gives the source disjunction on $[b,c)$.  Decreasing the returned
constant if needed
ensures $\kappa_{i+1}\le\kappa_i$.  The same argument applies independently
to every $r_{i+1}\in(0,r_i]$; choosing one admissible cutoff witness for each
radius produces the fixed selector $\mathfrak d_i$ required in the statement.
\end{proof}

\begin{theorem}[Bounded curvature at bounded distance]
\label{thm:bounded-curvature-at-bounded-distance}
\dcref{MT:ch10:10.2}
\source{morgan-tian}
\notready
For sufficiently small $\epsilon>0$ and every $A,C<\infty$, there are
$D_0=D_0(A,C,\epsilon)$ and $D=D(A,C,\epsilon)$ with the following
property.  Let $x$ lie in a three-dimensional Ricci flow with surgery whose
curvature is pinched toward positive.  Suppose every point $z$ at time at
most $\mathbf t(x)$ with $R(z)\ge4R(x)$ has a strong
$(C,\epsilon)$-canonical neighborhood.  If $R(x)\ge D_0$, then
\[
 R(y)\le D R(x)
 \qquad\text{for every }y\in
 B_{g(\mathbf t(x))}\bigl(x,A R(x)^{-1/2}\bigr).
\]
\end{theorem}

\begin{proof}
\uses{def:ricci-flow-with-surgery,def:curvature-pinched-toward-positive,def:strong-canonical-neighborhood-control,def:epsilon-neck-and-scale,lem:epsilon-neck-metric-control,prop:overlapping-epsilon-necks,lem:neck-chain-product,lem:nonflat-terminal-cone-exclusion}
Rescale so that $R(x)=1$ and suppose that points at bounded distance have
unbounded scalar curvature.  Along a minimizing geodesic from $x$, choose
the first points at which successive curvature levels are reached.  Above
level four, the canonical-neighborhood hypothesis applies.  A cap or compact
component at one of these first points has controlled diameter and curvature
ratio, so it cannot contain both the preceding controlled point and a point
at the next arbitrarily larger level.  The intervening regions are therefore
$\epsilon$-necks.

The metric and overlap estimates for necks assemble the relevant part of the
geodesic into a balanced chain.  By \ref{lem:neck-chain-product}, this chain
is a product tube.  Rescale at points whose neck radii tend to zero.  The
axial distance estimates and curvature pinching give a pointed incomplete
limit with nonnegative curvature, bounded curvature at one end, and
unbounded curvature at the other.  Blowing down the unbounded end, triangle
comparison shows that distances converge to
\[
 d_C((r,u),(s,v))^2=r^2+s^2-2rs\cos d_\Sigma(u,v),
\]
where the space of directions $\Sigma$ contains at least two points; the
limit is a genuine metric cone and not a ray.

The neck curvature derivative estimates provide uniform backward parabolic
neighborhoods on every compact subset away from the cone vertex.  Smooth
compactness therefore realizes the punctured cone as an open subset of a
nonnegatively curved Ricci flow at a positive time.  By
  \ref{lem:nonflat-terminal-cone-exclusion}, applied after translating the
 time interval so that the smooth cone slice is the terminal slice, rules out
 the nonflat cone.  The cone produced by the blow-up has scalar curvature one
 at its basepoint, so it cannot be the flat alternative.  This contradicts
 their construction.  The contradiction gives a uniform $D$;
choosing $D_0$ large enough to make the pinching and neck estimates uniform
proves the theorem.
\end{proof}

\begin{lemma}[Bounded curvature of a blow-up limit slice]
\label{lem:blowup-limit-slice-bounded-curvature}
\dcref{MT:ch11:11.7}
\source{morgan-tian}
\uses{thm:curvature-null-splitting}
\notready
Let $(M,g,x)$ be a complete, connected, nonflat, three-dimensional
Riemannian manifold with nonnegative curvature and $R_g(x)=1$.  Assume that
for every $A<\infty$ there is a compatible smooth pointed Ricci-flow limit on
$B_g(x,A)\times[-\tau_A,0]$, with $\tau_A>0$, whose zero-time slice is
$g|_{B_g(x,A)}$ and whose slices have nonnegative curvature.  Assume also
that every point $q$ with $R_g(q)>4$ has a strong $(2C,2\epsilon)$-canonical
neighborhood in the zero-time slice.  For sufficiently small $\epsilon$, the
curvature of $(M,g)$ is bounded on the whole slice:
\[
 \sup_{q\in M}|\operatorname{Rm}_g(q)|<\infty.
\]
The same conclusion applies to every time slice of a smooth pointed limit
flow whenever the stated limit-canonical-neighborhood hypothesis is supplied
at that time.
\end{lemma}

\begin{proof}
The limit-canonical-neighborhood premise is part of the Claim~11.7
interface; it is important that the constants are enlarged from
$(C,\epsilon)$ to $(2C,2\epsilon)$ before the claim is applied.

If the curvature operator has a zero eigenvalue at some point, the strong
maximum-principle and nullity argument in
\ref{thm:curvature-null-splitting}, applied on an exhaustion by bounded
parabolic neighborhoods, gives a one- or two-sheeted cover which is a product
$\Sigma^2\times\mathbb R$ with $\Sigma$ of positive curvature.  The complete
slice clause in that theorem makes the product global.  The surface factor is
a two-dimensional nonflat bounded-curvature noncollapsed ancient factor, so
the two-dimensional roundness clause of the same source contract makes it a
compact round surface (the lifted canonical neighborhood supplies the
required noncollapsing).  Hence its curvature is bounded, and the same is
true after taking the finite quotient.

If the curvature operator is everywhere positive and $M$ is compact, boundedness
is immediate.  Otherwise $M$ is noncompact.  A high-curvature canonical
neighborhood in a noncompact component can only be a strong neck or a cap with
a noncompact neck end; in either case it contains a $2\epsilon$-neck whose
scale is comparable to $R^{-1/2}$.  Unbounded curvature would therefore
produce $2\epsilon$-necks with scales tending to zero, contradicting the
no-arbitrarily-small-neck clause of the same source contract.  In dimension
three a
nonnegative curvature operator controls all components of $\operatorname{Rm}$
by the scalar curvature, so this gives the asserted global bound.
This is the three-case argument of Morgan--Tian Claim 11.7, with the
null-splitting and no-small-neck interfaces exposed above.
\end{proof}

\begin{lemma}[Closure of the noncollapsing predicate under pointed limits]
\label{lem:noncollapsing-predicate-pointed-limit}
\dcref{MT:ch11:11.1,MT:ch14:14.14,MT:ch14:14.15}
\source{morgan-tian}
\uses{def:surgery-flow-noncollapsing,thm:hamilton-compactness-generalized-flows}
Let smooth pointed parabolic regions converge in $C^\infty$ on compact
subsets, and suppose that every approximating admissible ball of scale
$\rho<\epsilon$ satisfies the uniform source noncollapsing predicate with
constant $\kappa$.  Then every admissible ball in the limit satisfies the
same volume-or-positive-component disjunction.  If the positive-component
alternative is absent on the limit region, the limit is ordinarily
$\kappa$-noncollapsed below $\epsilon$.
\end{lemma}

\begin{proof}
On a compactly contained unscathed parabolic ball, smooth convergence gives
convergence of the metrics, curvature bound, and ball volumes.  A strict
failure of the limiting volume inequality would therefore give the same
failure for the approximating balls, contradicting their uniform predicate.
If the limiting centre lies in a positive-sectional-curvature component, the
source disjunction records that alternative; otherwise the volume inequality
passes to the limit.  This is the volume-closure step used with Hamilton
compactness in Morgan--Tian Chapter~11, and no volume estimate is asserted
across a surgery boundary.
\end{proof}

\begin{theorem}[Pointed compactness for controlled flows]
\label{thm:pointed-controlled-flow-compactness}
\dcref{MT:ch11:11.1,MT:ch11:11.8,MT:ch11:11.10}
\source{morgan-tian}
\notready
Let $(\mathcal M_j,G_j,x_j,t_j)$ be based controlled three-dimensional Ricci
flows and put $Q_j=R_{G_j(t_j)}(x_j)\to\infty$.  Write
$\widetilde G_j(s)=Q_jG_j(t_j+Q_j^{-1}s)$, so that the base time is zero
and $R_{\widetilde G_j(0)}(x_j)=1$.  Assume curvature is pinched toward positive, and the rescaled
flows satisfy the uniform source predicate
\ref{def:surgery-flow-noncollapsing} on every fixed scale for $s<0$.  Every
point at nonpositive time with rescaled scalar curvature at least four has a
strong $(C,\epsilon)$-canonical neighborhood.  Suppose that for some
$0<T_0\le\infty$, every fixed
$B(x_j,0,A)\times(-T,0]$, with $A<\infty$ and $T<T_0$, eventually has a
compatible spacetime embedding.  In addition, assume that there are fixed
$\eta>0$ and $v_0>0$ such that, for all sufficiently large $j$,
\[
 \operatorname{Vol}_{\widetilde G_j(0)}
 B_{\widetilde G_j(0)}(x_j,\eta)\ge v_0.
\]
This is the volume-branch input required by Hamilton compactness; in the
first-failure application it is supplied by the ordinary noncollapsing bridge
and is not inferred from the positive-component exception in the surgery
noncollapsing predicate.  Then a subsequence converges smoothly to a complete
pointed Ricci flow on $(-T_0,0]$ with nonnegative curvature, locally bounded
in time, and the source noncollapsing disjunction passes to the limit.  If
$T_0=\infty$ and the positive-sectional-curvature alternative is absent on
every limit slice, the limit is a $\kappa$-solution; without that absence
hypothesis the conclusion is only the displayed disjunctive predicate.
\end{theorem}

\begin{proof}
\uses{def:surgery-flow-noncollapsing,def:controlled-ricci-flow-with-surgery,def:curvature-pinched-toward-positive,thm:hamilton-compactness-generalized-flows,lem:noncollapsing-predicate-pointed-limit,thm:local-ricci-flow-derivative-estimates,thm:bounded-curvature-at-bounded-distance,thm:curvature-null-splitting,lem:blowup-limit-slice-bounded-curvature,lem:finite-interval-harnack-distance-control}
On each fixed zero-time ball,
\ref{thm:bounded-curvature-at-bounded-distance} gives a scalar-curvature
bound.  Pinching converts the bound to a full curvature bound, and the
local derivative estimates propagate it for a short backward time depending
on the radius.  Together with the compatible embeddings and the displayed
base-ball volume estimate, \ref{thm:hamilton-compactness-generalized-flows}
therefore gives a compatible partial limit flow on
$B_{g_\infty(0)}(x_\infty,A)\times[-\tau_A,0]$ for some $\tau_A>0$.
Diagonalizing over integer $A$ gives the local-flow family required by
\ref{lem:blowup-limit-slice-bounded-curvature}; pointed Riemannian
compactness makes the zero-time slice complete.  The strict
canonical-neighborhood inequalities transfer to the limit with constants
enlarged to $(2C,2\epsilon)$, so the source Claim~11.7 interface supplies one
global curvature bound, not merely a bound depending on the radius of a based
ball.

With this global bound in hand, the local derivative estimates and the
flow-line lifetime hypothesis give a backward interval independent of $A$.
The displayed base-ball estimate gives the volume hypothesis in
\ref{thm:hamilton-compactness-generalized-flows}.  Diagonalizing over $A$
and compact backward-time intervals gives a smooth limit complete at time
zero.
The negative part of the curvature operator vanishes after division by
$Q_j$, hence the limit has nonnegative curvature; the final clause of the
compactness theorem now makes every earlier slice complete.

Apply \ref{lem:noncollapsing-predicate-pointed-limit} on each compact
unscathed parabolic region.  It passes the source volume-or-positive-component
disjunction to the limit and records exactly where an ordinary volume bound is
available.

To continue the limit backward, let $-T$ be a finite maximal endpoint with
$T<T_0$.  The finite-interval Harnack estimate
\ref{lem:finite-interval-harnack-distance-control} gives uniform diameter
control for every compact connected subset of the limit slices.  If such a
subset had a uniformly bounded scalar-curvature minimum but unbounded
curvature as $t\downarrow-T$, lift a bounded-curvature point and a
high-curvature point to the approximating flows.  Their distance stays
uniformly bounded, while the high-curvature point has a canonical
neighborhood; this contradicts
\ref{thm:bounded-curvature-at-bounded-distance}.  Thus every compact
subset has a uniform curvature bound once it contains a point of uniformly
bounded scalar curvature.

If a limit slice has a curvature-null direction, the product-cover conclusion
of \ref{thm:curvature-null-splitting} and the compact-factor clause of the
source Claim~11.7 interface (after normalizing at a positive-curvature point)
already give a global bound; the scalar-zero clause of the splitting theorem
gives flatness in the remaining case.  For a noncompact positive-curvature
limit, the exterior neck-separation and no-small-neck clauses in the source
Claims~11.11--11.16 give the same bound: a sequence of high-curvature
neighborhoods would otherwise contain separating necks with scales tending to
zero.  The compact case is covered because the scalar
minimum is nondecreasing under the flow and is at most
$R(x_\infty,0)=1$; hence each time has a point of uniformly bounded
curvature, so the compact-subset argument applies to the whole slice.  This
is the compact/exterior argument of Morgan--Tian Claims 11.11--11.16 and
rules out a finite-endpoint curvature blow-up.

The resulting global scalar bound, pinching, and the local derivative
estimates give uniform curvature bounds on
$B(x_j,0,A)\times[-(T+\delta),0]$ for every fixed $A$, with one
$\delta>0$ independent of $A$.  Hamilton compactness then extends the
limit to $(-(T+\\delta),0]$.  This is the continuation step of Morgan--Tian
Proposition 11.10; it uses the explicit compact/exterior bounds above rather
than inferring a global bound from a fixed-radius exhaustion.  Since every
$T<T_0$ is extendible, the limit reaches $(-T_0,0]$.
The finite-$T_0$ conclusion is only local boundedness in time; a fixed-radius
exhaustion does not control curvature at spatial infinity.  For every fixed
  limit time and every point of scalar curvature greater than four, the smooth
embeddings carry the point to approximating points with the same strict
curvature inequality.  The source limit-stability clause therefore transfers
a $(2C,2\epsilon)$-canonical neighborhood to the limit slice.  Applying
\ref{lem:blowup-limit-slice-bounded-curvature} at each time slice gives global
bounded nonnegative curvature.  Completeness and nonflatness
($R(x_\infty,0)=1$) are already established.  When the positive-component
alternative is absent, the closure lemma supplies ordinary all-scale
noncollapsing and hence exactly the defining properties of a $\kappa$-solution.
If that alternative remains, the theorem deliberately retains it rather than
silently promoting the limit to a $\kappa$-solution.
\end{proof}
